% concatenated from RegOct2025_Final.tex
%% This is a LaTeX template for preparing papers for Publ. Inst. Math.; version January 2016
%% Please delete everything begining with %% (DOUBLE %).

% Submission number: 
\documentclass[a4paper]{amsart}

\usepackage[british,english]{babel} 
\usepackage{graphics,color,pgf}
\usepackage{epsfig}
\usepackage[ansinew]{inputenc}
\usepackage[all]{xy}
\usepackage{hyperref}
\usepackage{amssymb, amsmath,amsthm, mathtools, amscd}
\usepackage{mathrsfs}
\usepackage{stmaryrd}
\usepackage{cancel} %strike out
\usepackage{tikz}
\usepackage{tikz-cd}
\usetikzlibrary{matrix}
%
%\theoremstyle{plain}
% \newtheorem{thm}{Theorem}
% \newtheorem{prop}[thm]{Proposition}
% \newtheorem{lem}[thm]{Lemma}
% \newtheorem{cor}[thm]{Corollary}
% \newtheorem*{op*}{Open Problem} 
% \newtheorem*{frthmA*}{Theorem A} 
% \newtheorem*{frcorB*}{Corollary B} 
%  \newtheorem*{frthmAp*}{Theorem A'} 
%  \newtheorem*{claim}{Claim} 
%\theoremstyle{definition}
% \newtheorem{exm}[thm]{Example}
% \newtheorem{dfn}[thm]{Definition}
%  \newtheorem{conj}[thm]{Conjecture}
%\theoremstyle{remark}
% \newtheorem*{rem}{Remark}

\newtheorem{intro_thm}{Theorem}
\newtheorem{intro_cor}[intro_thm]{Corollary}
\newtheorem{intro_prop}[intro_thm]{Proposition}


\newtheorem{lem}{Lemma}[section]
\newtheorem{thm}[lem]{Theorem} 
\newtheorem{prop}[lem]{Proposition}
\newtheorem{cor}[lem]{Corollary}
\newtheorem{conj}[lem]{Conjecture}
\newtheorem{fact}[lem]{Fact}
\newtheorem{claim}[lem]{Claim}
\newtheorem{ques}[lem]{Question}

\theoremstyle{definition}
\newtheorem{dfn}[lem]{Definition}
\newtheorem{exm}[lem]{Example}

\theoremstyle{remark}
\newtheorem{rem}[lem]{Remark}

\newtheoremstyle{TheoremNum}
        {0.2 cm}{0.2 cm}              %%% space between body and thm
        {\itshape}                      %%% Thm body font
        {}                              %%% Indent amount (empty = no indent)
        {}                     %%% Thm head font
        {.}                             %%% Punctuation after thm head
        { }                             %%% Space after thm head
        {\thmname{\bfseries #1}\thmnote{ \bfseries #3}}%%% Thm head spec
    \theoremstyle{TheoremNum}
   
\newtheorem{rec_thm}{Theorem}
\newtheorem{rec_prop}[rec_thm]{Proposition}
\newtheorem{rec_cor}[rec_thm]{Corollary}



%%New Math commands
\newcommand{\R}{\mathbb{R}} 
\newcommand{\N}{\mathbb{N}}
\newcommand{\Z}{\mathbb{Z}} 
\newcommand{\p}{\mathcal{P}} %piece
\newcommand{\s}{\mathcal{S}} %generating set
\newcommand{\myrightleftarrows}[1]{\mathrel{\substack{\xrightarrow{#1} \\[-.5ex] \xleftarrow{\quad\quad\;}}}}
\newcommand{\myrightleftarrowss}[2]{\mathrel{\substack{\xrightarrow{#1} \\[-.5ex] \xleftarrow[#2]{\quad\;}}}}
\DeclareMathOperator{\Ker}{Ker}
\DeclareMathOperator{\Img}{Im}
\DeclareMathOperator{\id}{\mathrm{id}}
\DeclareMathOperator{\h}{\mathrm{H}} %homology
\DeclareMathOperator{\C}{\mathrm{C}} %cochains
\DeclareMathOperator{\QM}{\mathrm{QM}}
\DeclareMathOperator{\incl}{\textbf{incl}}
\DeclareMathOperator{\restr}{restr}
\DeclareMathOperator{\alt}{\mathrm{alt}}
\DeclareMathOperator{\nalt}{\mathrm{n-alt}}
\DeclareMathOperator{\opp}{\mathrm{opp}}


%% Please, do not change the following four lines:
\renewcommand{\le}{\leqslant}\renewcommand{\leq}{\leqslant}
\renewcommand{\ge}{\geqslant}\renewcommand{\geq}{\geqslant}
\renewcommand{\setminus}{\smallsetminus}
\setlength{\textwidth}{28cc} \setlength{\textheight}{50cc} 


\title[Continuous cochains on Furstenberg boundaries]{Continuous cochains on Furstenberg boundaries and injectivity of the comparison map}

\author[Michelle Bucher]{%\bfseries 
Michelle Bucher} 
\address{Universit\'e de Gen\`eve}
\email{Michelle.Bucher@unige.ch}

\author[Alessio Savini]{%\bfseries 
Alessio Savini}
\address{University of Milano Bicocca} %% Please write ful names, avoid initials
\email{alessio.savini@unimib.it}



\thanks{Supported by the Swiss National Science Foundation. 
\\
\indent Mathematics Subject Classification 2020: Primary 22E41, Secondary 57T10.
} %% optional
%\thanks{Communicated by ...} %% This will be filled in the journal office.

\begin{document}

\begin{abstract}
 In \cite{Monod} Monod proved that any continuous cohomology of a semisimple Lie group $G$ can be represented by a measurable  cocycle on the associated Furstenberg boundary, which we upgraded to an alternating  cocycle in \cite{BucSavAlt}. In the current paper we improve that result by showing that we can actually take a representing cocycle which is continuous on an explicit subset of generic tuples. We give an analogous result in the case of bounded cohomology. Finally, we exploit this characterization to prove the injectivity of the comparison map in degree $3$ for $\mathrm{Isom}^\circ(\mathbb{H}_{\mathbb{C}}^n)$, when $n \geq 2$, and in degree $4$ for $\mathrm{Isom}^\circ(\mathbb{H}^n_{\mathbb{R}})$, when $n \geq 2$. 
\end{abstract}

\maketitle
\section{Introduction}

Let $G$ be a connected semisimple Lie group. It is a deep result by Austin and Moore  \cite[Theorem A]{AM} that the real valued continuous cohomology of $G$ is isomorphic to its measurable variant, namely $H^*_c(G) \cong H^*_m(G)$. The analogous result for bounded cohomology also holds, that is $H^*_{c,b}(G) \cong H^*_{m,b}(G)$ \cite[Proposition 7.5.1]{Monodbook}. In the bounded setting, one can prove that the latter bounded group cohomology is realizable on the Furstenberg boundary in the following sense: Let $P<G$ be a minimal parabolic subgroup and define
$$H^q_{m,b}(G \curvearrowright G/P):=H^q(L^\infty((G/P)^{*+1})^G),$$
where the cocomplex is endowed with the homogeneous coboundary operator. Evaluation on a base point induces a map
$$ev:H^q_{m,b}(G \curvearrowright G/P)\longrightarrow H^q_{m,b}(G)$$
which can easily be shown not to depend on the base point and further, exploiting the amenability of $P$, to be an isomorphism \cite[Corollary 7.5.9]{Monodbook}.  The analogous statement in the unbounded setting was recently shown to fail in higher rank: Defining 
$$H^q_{m}(G \curvearrowright G/P):=H^q(L^0((G/P)^{*+1})^G),$$
Monod showed in \cite{Monod} that the evaluation  map
$$ev:H^q_{m}(G \curvearrowright G/P)\longrightarrow H^*_{m}(G)$$
is surjective with a kernel that can explicitly be determined in terms of some invariant subspaces of $H^*_c(A)$, where $A<P$ is a maximal split torus. As this kernel becomes smaller when considering the corresponding cocomplexes with the additional condition that the cochains are alternating \cite{BucSavAlt}, from now on we will only consider alternating cochains. This has no effect on the measurable or continuous (bounded) cohomology of $G$, but $H^q_{m,\mathrm{alt}}(G \curvearrowright G/P)$ is in general a strict direct summand of $H^q_{m}(G \curvearrowright G/P)$.

 

In this paper we investigate the validity of the analogue of Austin and Moore's result for the measurable cohomology of $G$ on $G/P$. It is certainly too optimistic to expect an isomorphism with the continuous cohomology obtained by considering $G$-invariant continuous cochains on $G/P$. By continuity, such functions would automatically be bounded, but we exhibited in \cite{BucSavExp} several examples of cohomology classes in $H^*_{m}(G \curvearrowright G/P)$ which cannot be represented by bounded cocycles. Moreover, one of the simplest example of cocycles on $G/P$, namely the hyperbolic volume $\mathrm{Vol}_n$, assigning to every $(n+1)$-tuple $(x_0,\ldots,x_n)$ of points in $\partial \mathbb{H}^n_{\mathbb{R}}$ the signed volume of the associated ideal simplex, cannot be neither continuous nor cohomologous to a continuous cocycle. Indeed, given the existence of loxodromic isometries, any continuous function on $(\partial\mathbb{H}^n_\mathbb{R})^{n+1}$ has to be constant, and a constant alternating cochain is trivial. Nevertheless, the hyperbolic volume is continuous on tuples of pairwise distinct points. We will see that this phenomenon actually holds in full generality: with the right notion of \emph{genericity}, we will prove that every measurable cocycle on $G/P$ is actually cohomologous to another representative which is continuous on a full measure dense subset of tuples. 

%Roughly speaking, a triple of points $x,y,z \in G/P$ is \emph{generic} if $z$ is opposite to every Weyl chamber of the unique flat determined by $x$ and $y$. A $(q+1)$-tuple $(x_0,\ldots,x_q) \in (G/P)^{q+1}$ is \emph{generic} if for any choice of distinct indices $0 \leq i,j,k \leq q$, the triple $(x_i,x_j,x_k)$ is generic. 

We postpone to Section \ref{sec barycenter} the precise Definition \ref{Def generic} and we denote by $(G/P)^{(q+1)}$ the subset of generic tuples. For now we record the following facts:
\begin{itemize}
\item  The subset $(G/P)^{(q+1)}$ is open and dense in $(G/P)^{q+1}$, thus it has full measure. 
\item  When $G$ has real rank equal to one, the subset of generic tuples is precisely the subset of pairwise distinct tuples. 
\end{itemize}

We define the \emph{continuous alternating cohomology} of the $G$-action on $G/P$ as the cohomology of the cocomplex
$$
H^q_{c,\alt}(G \curvearrowright G/P):=H^q(C_{c,\alt}((G/P)^{(*+1)})^G),
$$
where 
$$
C_{c,\alt}((G/P)^{(q+1)})^G:=\{ f : (G/P)^{(q+1)} \rightarrow \mathbb{R} \ | \ \textup{$f$ is continuous and alternating} \}^G
$$
is endowed with the homogeneous coboundary operator. 


\begin{intro_thm}\label{thm_continuous}
Let $G$ be a connected semisimple Lie group with finite center and $P<G$ be a minimal parabolic subgroup. The inclusion of cocomplexes 
$$
C_{c,\alt}((G/P)^{(q+1)})^G \longrightarrow L^0_{\alt}((G/P)^{q+1})^G 
$$
induces a surjection 
$$
\xymatrix{
H^q_{c,\alt}(G \curvearrowright G/P) \ar@{->>}[r]&  H^q_{m,\alt}(G \curvearrowright G/P).
}
$$
\end{intro_thm}

As a consequence any class in the continuous cohomology of $G$ can be actually realized as a $G$-invariant alternating cocycle on the boundary $G/P$ which is continuous on the subset of generic tuples. When $G$ has real rank equal to one this subset is the subset of pairwise distinct tuples. 

We do not know if the surjection of Theorem \ref{thm_continuous} is an isomorphism in general, as one could wish for when looking for an anologue of the result by Austin and Moore. The situation improves in the bounded setting. Defining the \emph{continuous bounded alternating cohomology} of the $G$-action on $G/P$ as 
$$
H^q_{c,b,\alt}(G \curvearrowright G/P):=H^q(C_{c,b,\alt}((G/P)^{(\ast+1)})^G),
$$
we obtain the desired isomorphism: 

\begin{intro_thm}\label{thm_continuous_bounded}
Let $G$ be a connected semisimple Lie group with finite center and $P<G$ be a minimal parabolic subgroup. The inclusion of cocomplexes 
$$
C_{c,b,\alt}((G/P)^{(q+1)})^G \longrightarrow L^\infty_{\alt}((G/P)^{q+1})^G 
$$
induces an isomorphism
$$H^q_{c,b,\alt}(G \curvearrowright G/P)\cong  H^q_{m,b,\alt}(G \curvearrowright G/P) .$$
\end{intro_thm}


\subsection*{Injectivity of the comparison map} The forgetful functor induces so called \emph{comparison maps} which naturally fit in the commutative diagram

\begin{equation}\label{diagram comparison intro}
\xymatrix{
 H^q_{c,b,\alt}(G \curvearrowright G/P) \ar[d]_{\mathrm{comp}^q_{G \curvearrowright G/P}} \ar@{->}[rr]^{ ev}_{\cong} &&H^q_{c,b}(G) \ar[d]^{\mathrm{comp}^q_G}\\
H^q_{c,\alt}(G \curvearrowright G/P)\ar@{->}[rr]^{ ev} && H^q_c(G).
}
\end{equation}

It is a mysterious open question whether the function $\mathrm{comp}_G^q$ is actually an isomorphism \cite[Problem A]{MonodICM}. In the current manuscript we focus our attention on the injectivity issue. The injectivity of $\mathrm{comp}^q_G$ is trivial in degrees $0$ and $1$. Injectivity is known to hold:
\begin{itemize}
\item In degree $2$ in full generality \cite{BM99}.
\item In degree $3$ for $\mathrm{SL}(2,\mathbb{R})$ \cite{BuMo2}, for $\mathrm{SL}(2,\mathbb{C})$ \cite{Bloch}, for $\mathrm{SL}(n,\mathbb{R})$ \cite{MonodSLn}, for $\mathrm{SL}(n,\mathbb{C})$ \cite{MonodSLn,BBI18}, for $\mathrm{Isom}^\circ(\mathbb{H}^n_{\mathbb{R}})$ \cite{Pieters}, for $\mathrm{SO}(n,\mathbb{C})$ and $\mathrm{Sp}(2n,\mathbb{C})$ \cite{DLC}, for products of groups of orientation preserving isometries of hyperbolic spaces \cite{BucSavExp}. 
\item In degree $4$ only for $\mathrm{SL}(2,\mathbb{R})$ \cite{HO}.
\end{itemize} 

In this paper we further establish injectivity in degree $3$ for $\mathrm{Isom}^\circ(\mathbb{H}^n_\mathbb{C})$ (Corollary \ref{cor complex}) and in degree $4$ for  $\mathrm{Isom}^\circ(\mathbb{H}^n_\mathbb{R})$ (Theorem \ref{thm real}). 


From Diagram (\ref{diagram comparison intro}) it is clear that $\mathrm{comp}^q_G$ is injective if and only if $\mathrm{comp}^q_{G \curvearrowright G/P}$ is injective and the kernel of $ev$ intersects the image of $\mathrm{comp}^q_{G \curvearrowright G/P}$ trivially. We will see that exploiting the proofs of Theorems \ref{thm_continuous} and \ref{thm_continuous_bounded} we can do a little better by restricting the latter condition on the  isomorphic copy of 

$$H^{q-1}_c(A)^{w_0}=\mathrm{Ker}(ev: H^q_{m,\mathrm{alt}}(G \curvearrowright G/P)\longrightarrow H^q_{m}(G)),$$
which naturally injects in the kernel of the evaluation map  
$$ev: H^q_{c,\mathrm{alt}}(G \curvearrowright G/P)\longrightarrow H^q_{c}(G).$$
More precisely, we establish the following criterion: 

\begin{intro_prop}\label{prop injectivity comparison}\label{criterion}
Let $G$ be a connected semisimple Lie group with finite center. Let $P$ be a miminal parabolic subgroup, $A<P$ a maximal split torus and $w_0$ a representative of the longest element in the Weyl group. The comparison map 
$$
\mathrm{comp}^q_G:H^q_{c,b}(G) \rightarrow H^q_c(G)
$$
is injective if and only if 
\begin{enumerate}
\item the kernel of the measurable evaluation map consists of unbounded classes in $H^q_{c,\alt}(G \curvearrowright G/P)$ , i.e. 
$$H^{q-1}_c(A)^{w_0}\cap \mathrm{comp}^q_{G \curvearrowright G/P}(H^q_{c,b,\alt}(G \curvearrowright G/P))=\{0\},$$ 

\item the comparison map 
$$
\mathrm{comp}^q_{G \curvearrowright G/P}: H^q_{c,b,\alt}(G \curvearrowright G/P) \longrightarrow H^q_{c,\alt}(G \curvearrowright G/P)
$$
is injective.
\end{enumerate}

\end{intro_prop} 


The first condition of this criterion is trivially satisfied for $q$ strictly greater than the real rank of $G$. It also always holds for even $q$ when $w_0$ acts as $-1$ on the Lie algebra of $A$, since in that case $H^{q-1}_c(A)^{w_0}=0$. It is our belief that the validity of the first condition is the easier part of this criterion. The difficulty of proving injectivity of the comparison map has thus shifted to the boundary. Exploiting continuity and transitivity properties of the action of $\mathrm{Isom}^\circ(\mathbb{H}^n_{\mathbb{R}})$ on its Furstenberg boundary we prove: 

\begin{intro_thm}\label{thm boundary deg 4}
Let $G=\mathrm{Isom}^\circ(\mathbb{H}^n_{\mathbb{R}})$, where $n \geq 2$. If $P<G$ is any parabolic subgroup, the map
$$
\xymatrix{
\mathrm{comp}^4_{G \curvearrowright G/P}:H^4_{c,b}(G \curvearrowright G/P) \ar@{^{(}->}[r] & H^4_c(G \curvearrowright G/P)
}
$$
is injective. 
\end{intro_thm}

As a consequence of Proposition \ref{prop injectivity comparison} and Theorem \ref{thm boundary deg 4} we obtain 

\begin{intro_thm}\label{thm group deg 4}\label{thm real}
Let $G=\mathrm{Isom}^\circ(\mathbb{H}^n_{\mathbb{R}})$, where $n \geq 2$. The comparison map 
$$
\xymatrix{
\mathrm{comp}^4_{G}:H^4_{c,b}(G) \ar[r]^{\hspace{10pt} \cong} & H^4_c(G)
}
$$
is an isomorphism. 
\end{intro_thm}

This was known previously only in the case of $n=2$ by a tour de force by Hartnick and Ott  \cite{HO} involving elaborate partial differential equations in a way we cannot claim to fully comprehend. Our proof appears to us as more elementary. Note that the surjectivity is well known. 


\subsection*{Configurations of triples of points and injectivity in degree $3$} We conclude with a curious interplay between an algebraic property of the action of the longest element $w_0$ and the topology (more precisely only the non-compactness) of the configuration space of triples of generic points:

\begin{intro_prop}\label{intro prop injective degree 3}
Let $G$ be a connected semisimple Lie group with finite center. Let $P<G$ be a minimal parabolic subgroup, $A<P$ a maximal split torus and $w_0$ a representative of the longest element in the Weyl group. \begin{enumerate}
\item If the quotient $G\backslash (G/P)^{(3)}$ is compact, then the comparison map
$$
\mathrm{comp}^2_{G \curvearrowright G/P}:H^2_{c,b}(G \curvearrowright G/P) \longrightarrow H^2_c(G \curvearrowright G/P)
$$
is surjective and 
$$
\mathrm{comp}^3_{G \curvearrowright G/P}:H^3_{c,b}(G \curvearrowright G/P) \longrightarrow H^3_c(G \curvearrowright G/P)
$$
is injective. 
\item If $w_0$ does not act as $-1$ on $\mathfrak{a}=\mathrm{Lie}(A)$, then $G\backslash (G/P)^{(3)}$ is not compact.
\end{enumerate}
\end{intro_prop}

\begin{proof} \begin{enumerate} \item A $G$-invariant continuous cochain $f:(G/P)^{(3)}\rightarrow \mathbb{R}$ corresponds to a continuous map $F:G\backslash (G/P)^{(3)}\rightarrow \mathbb{R}$. If $G\backslash (G/P)^{(3)}$ is compact, such a map is automatically bounded, and thus in this case any $G$-invariant continuous cochain $f:(G/P)^{(3)}\rightarrow \mathbb{R}$ is bounded. This immediately proves surjectivity in degree $2$. For the injectivity in degree $3$, suppose that a bounded $G$-invariant continuous cocycle $b: (G/P)^{(4)}\rightarrow \mathbb{R}$ vanishes in $H^3_c(G \curvearrowright G/P)
$. This means that there exists a $G$-invariant continuous cochain $f:(G/P)^{(3)}\rightarrow \mathbb{R}$ such that $\delta f=b$. But if $G\backslash (G/P)^{(3)}$ is compact, the cochain $f$ is bounded. 
\item We will use injectivity of the comparison map in degree $2$  \cite{BM99}, which by Proposition \ref{criterion} in particular implies that the kernel of the evaluation map intersects the image of the comparison map 
$$
\mathrm{comp}^2_{G \curvearrowright G/P}:H^2_{c,b}(G \curvearrowright G/P) \longrightarrow H^2_c(G \curvearrowright G/P)
$$
trivially. Now $w_0$ does not act as $-1$ if and only if $H^1_c(A)^{w_0}\neq 0$. In particular the kernel of the evaluation map is nontrivial and as a consequence the latter comparison map is not surjective. By the first item of the proposition this implies that the quotient $G\backslash (G/P)^{(3)}$ is not compact. 
\end{enumerate}
\end{proof}

For example, for $G=\mathrm{Isom}^\circ(\mathbb{H}^2_\mathbb{R})$, there are two orbits of generic triple of points in the Furstenberg boundary $\partial \mathbb{H}^2_\mathbb{R}$ so that the configuration space of distinct triples of points consists of $2$ points and is hence compact. For $G=\mathrm{Isom}^\circ(\mathbb{H}^n_\mathbb{R})$, with $n\geq 3$, there is only one orbit. For $G=\mathrm{Isom}^\circ(\mathbb{H}^n_\mathbb{C})$, with $n\geq 2$, the Cartan angular invariant gives a homeomorphism between the configuration space of distinct triples of points in $\partial \mathbb{H}^n_\mathbb{C}$ and the closed interval $[-\pi/2,\pi/2]$. In contrast, for $G=\mathrm{SL}(3,\mathbb{R})$, the Furstenberg boundary is given by the space of maximal flags in $\mathbb{R}^3$ and the triple ratio of generic triples of maximal flags induces a homeomorphism between the configuration space of triples and $\mathbb{R}\setminus \{0,1\}$ which  is indeed not compact. 

Note that the converse of the second item of Proposition \ref{intro prop injective degree 3} does not hold. For example for $G=\mathrm{Sp}(4,\mathbb{R})$ the longest element $w_0$ acts on $\mathfrak{a}$ as $-1$ but the quotient $G\backslash (G/P)^{(3)}$ is not compact. 

We conclude the introduction with a direct application of Proposition \ref{intro prop injective degree 3}: 

\begin{intro_cor}\label{cor complex}
Let $G=\mathrm{Isom}^\circ(\mathbb{H}^n_{\mathbb{C}})$, where $n \geq 2$. The comparison map
$$
\xymatrix{
\mathrm{comp}^3_{G}:H^3_{c,b}(G) \ar@{^{(}->}[r] & H^3_c(G).
}
$$
is injective. 
\end{intro_cor}

\begin{proof} We only need to verify the two items of Proposition \ref{criterion}: The first one holds since $A\cong \mathbb{R}$ so that $H^2_c(A)^{w_0}=0$. The second holds since, as noticed above, the space of configuration of triples of points is homeomorphic to a closed interval, so that the comparison map 
$$
\xymatrix{
\mathrm{comp}^3_{G \curvearrowright G/P}:H^3_{c,b}(G \curvearrowright G/P) \ar@{^{(}->}[r] & H^3_c(G \curvearrowright G/P).
}
$$
is injective by the first item of Proposition \ref{intro prop injective degree 3}. 
\end{proof}



%\begin{dfn}\label{dfn_weyl_opposition}
%Given $n \in N$, we say that the triple $(P,w_0P,nw_0P)$ satisfies the condition of \emph{Weyl opposition} if for any element $w \in W$ in the Weyl group we have that $nwP$ is opposite to $w_0P$. A triple $(g_0P,g_1P,g_2P)$ is Weyl opposite if it lies in the same $G$-orbit of a triple $(P,w_0P,nw_0P)$ satisfying Weyl opposition.
%
%For $q \geq 3$, a $q$-tuple $(g_1P,\ldots,g_qP)$ is \emph{Weyl opposite} if any extracted triple $(g_iP,g_jP,g_kP)$ is Weyl opposite, for $i<j<k$. We denote by $(G/P)^{(q)}$ the subset of Weyl opposite $q$-tuples. 
%\end{dfn}
%
%The notion of Weyl opposition is clearly empty if we deal with only one point and we identify it with the usual opposition when we refer to pairs of points. When $G$ is a rank one Lie group, the Weyl group  has order two and Weyl opposite tuples boil down to \emph{pairwise opposite tuples}. In the higher rank case, Weyl opposition is actually a stronger condition than pairwise opposition.
%
%\begin{lem}\label{lem_weyl_full}
%\end{proof}
%
%\begin{rem}
%In the proof of Theorem \ref{thm barycenter} we are going to show that $(G/P)^{(3)}$ is diffeomorphic to the quotient $(G \times N_{\opp})/MA$, where $MA$ acts on $G$ by right multiplication and on $N_{\opp}$ by conjugation. 
%\end{rem}
%
%Our aim is to use the characterization of Weyl opposite triples to obtain suitable projections on the maximal split torus $A$. We start noticing that, if $n \in N_{\opp}$, then $nwP$ is opposite to $w_0P$ for every $w \in W$. By multiplying both points by $w_0$, we have that $w_0nwP$ is opposite to $P$ (since $w_0^2 \in M <P$). As a consequence of Equation \eqref{eq:N:parametrization}, there must exist a unique
%$$
%\iota_w(n) \in N
%$$
%such that 
%\begin{equation}\label{eq iota} 
%w_0nwP=\iota_w(n)w_0P. 
%\end{equation}
%The above notation is not precise, since $\iota_w(n)$ depends also on the choice of the representative $w_0$ of the longest element in the Weyl group. In virtue of the previous observation, one could ask how the function $\iota_w$ may vary if we change the representative $w_0$. 
%
%\begin{lem}\label{iotaw depend w0}
%Let $w_0$ be a representative of the longest element in the Weyl group and consider another representative of the form $w_0\rho$, where $\rho \in M$. If we denote by $\iota_w^{w_0}$ the function depending on $w_0$ and by $\iota_w^{w_0\rho}$ the function dependening on $w_0\rho$, we have that 
%$$
%\iota_w^{w_0\rho}(n)=\rho_{w_0}\iota_w^{w_0}(n)\rho_{w_0}^{-1},
%$$
%for all $n \in N_{\opp}$, where $\rho_{w_0}=w_0 \rho w_0^{-1}$. 
%\end{lem}
%
%\begin{proof}
%In this context, we have that
%$$
%nwP=\rho^{-1}w_0^{-1}\iota_w^{w_0\rho}(n)w_0P=w_0^{-1}\rho_{w_0}^{-1}\iota_w^{w_0\rho}(n)\rho_{w_0}w_0P,
%$$
%where we exploited the fact that $w_0$ normalizes $M$. Notice that the product $\rho_{w_0}^{-1}\iota_w^{w_0\rho}(n)\rho_{w_0}$ is an element of $N$, since the latter is preserved by $M$-conjugation. 
%
%At the same time, we must have
%$$
%nwP=w_0^{-1}\iota_w^{w_0}(n)w_0P,
%$$ 
%thus, comparing the two equations, Equation \eqref{eq:N:parametrization} implies that
%$$
%\iota_w^{w_0\rho}(n)=\rho_{w_0}\iota_w^{w_0}(n)\rho_{w_0}^{-1}. 
%$$
%and the statement follows. 
%\end{proof}
%
%\begin{exm}\label{exm Hn iota}
%Consider $G=\mathrm{Isom}^+(\mathbb{H}^n)$. If we realize the boundary at infinity of $\mathbb{H}^n$ as $\mathbb{R}^{n-1} \cup \{ \infty \}$, we can define $P=\mathrm{Stab}_G(\infty)$ and take as maximal split torus $A$ the subgroup of homotheties $\{ a_\lambda: x \mapsto \lambda x \ | \ \lambda > 0\}$. In this context, we can identify the unipotent radical $N$ of $P$ as the subgroup of translations, namely 
%$$
%\mathbb{R}^{n-1} \longrightarrow N, \ \ v \mapsto n_v(x):=x+v.
%$$
%
%We fix as representative of the longest element in the Weyl group the following involution
%$$
%w_0:\mathbb{R}^n \rightarrow \mathbb{R}^n, \ \ w_0(x):=\frac{x-2\langle e_1,x\rangle e_1}{\lVert x \rVert^2},
%$$
%where $e_1=(1,0,\ldots,0)$ is the first vector of the canonical basis of $\mathbb{R}^n$, $\langle \cdot \ , \ \cdot\rangle$ denotes the standard scalar product and $\lVert \ \cdot \ \rVert$ the associated norm. In this way, one can check that the map $i_{w_0}$ is given by 
%$$
%i_{w_0}(n_v)=n_{w_0(v)}.
%$$
%\end{exm}
%
%\begin{exm}\label{exm SL3 iota}
%Given $G=\mathrm{SL}(3,\mathbb{R})$, we can choose as $P$ the subgroup of upper triangular matrices and we define the maximal split torus as
%$$
%A=\{ \mathrm{diag}(a_{11},a_{22},a_{33}) \ | \ a_{11}a_{22}a_{33}=1, \ a_{ii} >0 \textup{\ \ for $i=1,2,3$} \}. 
%$$
%The unipotent radical $N$ of $P$ is given by the subgroup of unipotent matrices. By fixing as $w_0$ the matrix
%$$
%w_0=
%\left(
%\begin{array}{ccc}
%0 & 0 & -1 \\
%0 & 1 & 0  \\
%1 & 0 & 0 
%\end{array}
%\right),
%$$
%we obtain that 
%$$
%i_{w_0}
%\left(
%\begin{array}{ccc}
%1 & x & y \\
%0 & 1 & z \\
%0 & 0 & 1 \\
%\end{array}
%\right)=
%\left(
%\begin{array}{ccc}
%1 & \frac{x}{y-xz} & -\frac{1}{y}\\
%0 & 1 & \frac{z}{y} \\
%0 & 0 & 1\\ 
%\end{array}
%\right).
%$$
%In particular, if $n \in N_{opp}$, we must have $y\neq 0$ and $xz-y \neq 0$. 
%\end{exm}
%
%\begin{lem}
%The subset $N_{\opp} \subset N$ is invariant by  $MA$-conjugation. Additionally, for any $w \in W$, we have the following relations:
%\begin{equation}\label{eq A conjugation}
%\iota_w(ana^{-1})=a_{w_0}\iota_w(n) a_{w_0}^{-1},
%\end{equation}
%\begin{equation}\label{eq M conjugation}
%\iota_w(\rho n \rho^{-1})=\rho_{w_0} \iota_w(n) \rho_{w_0}^{-1}, 
%\end{equation}
%where $a_{w_0}=w_0aw_0^{-1}$ and $\rho_{w_0}=w_0\rho w_0^{-1}$ for any $a \in A, \rho \in M$. 
%\end{lem}
%
%\begin{proof}
%We are going to focus our attention only on the $A$-action and we leave to the reader the case of $M$. Consider an element $a \in A$ and define $a_{w_0}=w_0aw_0^{-1}$. If we start with the equation 
%$$
%nwP=w_0^{-1}\iota_w(n)w_0P,
%$$
%we can multiply both sides by $a$, obtaining
%\begin{equation}\label{eq A action}
%anwP=aw_0^{-1}\iota_w(n)w_0P. 
%\end{equation}
%The left-hand side of Equation \eqref{eq A action} can be written as
%$$
%anwP=ana^{-1}awP=ana^{-1}wP,
%$$
%where the last equality follows by the fact that the conjugation by $w$ preserves $A$, which is a subgroup of $P$. Notice that $ana^{-1} \in N$, since $N$ is normalized by $A$. 
%
%
%
%If we now consider the right-hand side of Equation \eqref{eq A action}, we get that
%$$
%aw_0^{-1}\iota_w(n)w_0P=w_0^{-1}a_{w_0}\iota_w(n)a_{w_0}^{-1}a_{w_0}w_0P=w_0^{-1}a_{w_0}\iota_w(n)a_{w_0}^{-1}w_0P,
%$$
%where the first equation follows be the definition of $a_{w_0}$ and the second one follows by the fact the conjugation by $w_0$ preserves $A$. The element $a_{w_0}i_w(n)a_{w_0}^{-1}$ is still an element of $N$, since the latter is preserved by $A$-conjugation. To sum up, we obtain the following equation
%$$
%w_0ana^{-1}wP=a_{w_0}i_w(n)a_{w_0}^{-1}w_0P. 
%$$
%By the defining Equation \eqref{eq iota}, the uniqueness of $\iota_w(ana^{-1})$ implies that 
%$$
%\iota_w(ana^{-1})=a_{w_0}\iota_w(n)a_{w_0}^{-1},
%$$
%and the claim is proved. 
%\end{proof}
%
%The Iwasawa decomposition \cite[Theorem IX.1.3]{Helgason} implies that any element $g \in G$ can be written as a product $ank$, where $a \in A, n \in N, k \in K$. In particular, we are allowed to define the $A$-\emph{projection} as follows
%$$
%\pi_A:G \longrightarrow A, \ \ \ \pi_A(g)=a. 
%$$
%By the uniqueness of the decomposition, we have that 
%\begin{equation}\label{M invariance projection}
%\pi_A(\rho g)=\pi_A(g), 
%\end{equation}
%for any $\rho \in M$. In fact, an easy computation shows that
%$$
%\rho g=\rho a n k = a \rho n k= a n' \rho k,
%$$
%where we exploited the fact that $\rho$ centralizes $A$ and normalizes $N$. The uniqueness of the decomposition allows us to conclude. 
%
%\begin{dfn}\label{def psi}
%For any element $w \in W$, we define the $w$-\emph{projection} on $A$ as
%$$
%\psi_w:N_{\opp} \rightarrow A, \ \ \  \psi_w(n):=\pi_A(w^{-1}n^{-1}w_0^{-1}\iota_w(n)w_0),
%$$
%and similarly we define
%$$
%\Psi:N_{\opp} \rightarrow A, \ \ \ \Psi(n):=\prod_{w \in W} (\psi_w(n))^{-\frac{1}{|W|}},
%$$
%where $|W|$ denotes the cardinality of the Weyl group. 
%\end{dfn}
%
%It is worth noticing that the exponent $1/|W|$ makes sense, since we are simply looking at the corresponding element in the Lie algebra $\mathrm{Lie}(A)=\mathfrak{a}$ 
%via the logarithm, then we multiply by $1/|W|$ and we exponentiate back. In virtue of Equation \eqref{M invariance projection} the function $\psi_w$ does not depend on the choice of the representative of $w$, whereas it may depend \emph{a priori} on the choice of the representative of $w_0$. 
%
%\begin{lem}\label{lem MA psiw}
%For any element $w \in W$, the function $\psi_w$ does not depend on the choice of the representative of $w_0$. Additionally we have that
%\begin{equation}\label{eq psiw M invariance}
%\psi_w(\rho n \rho^{-1})=\psi_w(n),
%\end{equation}
%\begin{equation}\label{eq psiw A invariance}
%\psi_w(a n a^{-1})=(w^{-1}aw)a^{-1}\psi_w(n), 
%\end{equation}
%for any $a \in A, \rho \in M$. 
%\end{lem}
%
%\begin{proof}
%As in the previous lemma, we set
%$$
%a_{w_0}=w_0aw_0^{-1}, \ \ \ \rho_{w_0}=w_0\rho w_0^{-1}. 
%$$
%We start showing that $\psi_w$ does not depend on the particular representative of $w_0$. We consider the element $w_0 \rho$, where $\rho \in M$ and we use the same notation $\iota_w^{w_0\rho}$ adopted in Lemma \ref{iotaw depend w0}. 
%\begin{align*}
%\pi_A(w^{-1}n^{-1}\rho^{-1}w_0^{-1}\iota_w^{w_0\rho}(n)w_0\rho)&=\pi_A(w^{-1}n^{-1}w_0^{-1}\rho_{w_0}^{-1}\iota_w^{w_0\rho}(n)\rho_{w_0}w_0)\\
%&=\pi_A(w^{-1}n^{-1}w_0^{-1}\iota_w(n)w_0)\\
%&=\psi_w(n),
%\end{align*}
%where we exploited Lemma \ref{iotaw depend w0} to pass from the first line to the second one. Thus $\psi_w$ is independent of the choice of the representative $w_0$. 
%
%If we consider $\rho \in M$, we have that
%\begin{align*}
%\psi_w(\rho n \rho^{-1})&=\pi_A(w^{-1}\rho n^{-1} \rho^{-1} w_0^{-1}\iota_w(\rho n \rho^{-1}) w_0)\\
%&=\pi_A(w^{-1}\rho w w^{-1} n^{-1} \rho^{-1} w_0^{-1} \rho_{w_0} \iota_w(n) \rho_{w_0}^{-1} w_0)\\
%&=\pi_A(w^{-1} n^{-1} w_0^{-1} \iota_w(n) w_0 \rho^{-1})\\
%&=\pi_A(w^{-1} n^{-1} w_0^{-1} \iota_w(n) w_0)=\psi_w(n), 
%\end{align*}
%where we moved from the first line to the second one using Equation \eqref{eq M conjugation}, we passed from the second line to the third one exploiting Equation \eqref{M invariance projection} and the definition of $\rho_{w_0}$ and we concluded thanks to the fact that the evalutation of $\pi_A$ on an element of $M$ on the right is trivial.
%
%Consider now an element $a \in A$. We can compute the following projection
%\begin{align*}
%\psi_w(ana^{-1})&=\pi_A(w^{-1}an^{-1}a^{-1}w_0^{-1}\iota_w(ana^{-1})w_0)\\
%&=\pi_A(w^{-1}aw w^{-1} n^{-1}a^{-1}w_0^{-1}a_{w_0}\iota_w(n)a_{w_0}^{-1}w_0)\\
%&=(waw^{-1})\pi_A(w^{-1}n^{-1}w_0^{-1}\iota_w(n)w_0a^{-1})\\
%&=(waw^{-1})a^{-1}\psi_w(n),
%\end{align*}
%where we passed from the first line to the second one using Equation \eqref{eq A conjugation}, we moved from the second line to the third one exploiting the definition of $a_{w_0}$ and the fact that $waw^{-1}$ is in $A$, and we concluded by observing that $w^{-1}n^{-1}w_0^{-1}\iota_w(n)w_0 \in P$. In fact, for every $p \in P$, it holds 
%$$
%\pi_A(p a)=\pi_A(p)a,
%$$
%because $a$ commutes with $M$ and normalizes $N$. This concludes the proof. 
%\end{proof}
%
%\begin{exm}
%Let $G=\mathrm{SL}(3,\mathbb{R})$. We fix the same subgroups $P$ and $A$ as in Example \eqref{exm SL3 iota}. If we consider the representative of $w_0$ given by 
%$$
%w_0=
%\left(
%\begin{array}{ccc}
%0 & 0 & -1 \\
%0 & 1 & 0  \\
%1 & 0 & 0  \\
%\end{array}
%\right),
%$$
%the function $\psi_{w_0}$ we obtain is given by
%$$
%\psi_{w_0}
%\left(
%\begin{array}{ccc}
%1 & x & y\\
%0 & 1 & z\\
%0 & 0 & 1\\
%\end{array}
%\right)=
%\left(
%\begin{array}{ccc}
%1/y & 0 & 0\\
%0 & -\frac{y}{xz-y} &0\\
%0 & 0 & y-zx
%\end{array}
%\right).
%$$
%\end{exm}
%
%
%\begin{prop}\label{prop Psi properties}
%The function $\Psi:N_{\opp} \rightarrow A$ satisfies the following properties 
%\begin{equation}\label{eq Psi M invariant}
%\Psi(\rho n \rho^{-1})=\Psi(n),
%\end{equation}
%\begin{equation}\label{eq Psi A equivariant}
%\Psi(a n a^{-1})=a \Psi(n),
%\end{equation}
%for all $\rho \in M,a \in A$. 
%\end{prop}
%
%\begin{proof}
%The fact that $\Psi$ is invariant under $M$-conjugation is an immediate consequence of Equation \eqref{eq psiw M invariance}. We are left to show the $A$-equivariance of $\Psi$. 
%
%\begin{align*}
%\Psi(ana^{-1})&=\left( \prod_{w \in W} \psi_w(ana^{-1})\right)^{-\frac{1}{|W|}}\\
%&=\left( \prod_{w \in W} (waw^{-1})a^{-1}\psi_w(n) \right)^{-\frac{1}{|W|}}\\
%&=(a^{-|W|})^{-\frac{1}{|W|}} \left( \prod_{w \in W} \psi_w(n) \right)^{-\frac{1}{|W|}}\\
%&=a \Psi(n),
%\end{align*}
%where we exploited the fact that $\prod_{w \in W} (waw^{-1})=e$ to move from the second line to the third one and we conclude using Equation \eqref{eq psiw A invariance}. 
%\end{proof}
%
%We conclude the section by showing the existence of a continuous equivariant map from Weyl opposite triples to the associated symmetric space. 
%
%\begin{rec_thm}[\ref{thm barycenter}]
%Let $G$ be a connected semisimple Lie group. Let $P$ be a minimal parabolic subgroup and denote by $K$ a maximal compact subgroup. Then, there exists a $G$-equivariant continuous map 
%$$
%\Phi: (G/P)^{(3)} \rightarrow G/K,
%$$
%from the set of Weyl opposite triples to the symmetric space associated to $G$. 
%\end{rec_thm}
%
%\begin{proof}
%Recall that we have a continuous surjection given by
%$$
%\pi:G \times N_{\opp} \longrightarrow (G/P)^{(3)}, \ \ \pi(g,n):=(gP,gw_0P,gnw_0P). 
%$$
%We define the function
%$$
%\varphi:G \times N_{\opp} \longrightarrow G/K, \ \ \varphi(g,n):=g\Psi(n)K,
%$$
%which is continuous because the $G$-action is continuous and the function $\Psi$ is continuous (being the composition of continuous functions). Our goal is to use the universal property of quotients to get the following commutative diagram
%\begin{equation}\label{diagram section}
%\xymatrix{
%G \times N_{\opp} \ar[rr]^{\varphi} \ar[dr]_{\pi} && G/K\\
%& (G/P)^{(3)}. \ar@{.>}[ur]_{\Phi}
%}
%\end{equation}
%
%In order to do that, it is sufficient to check that we can factor $\varphi$ through $\pi$, namely that $\varphi(n,g)=\varphi(g',n')$ whenever $\pi(g,n)=\pi(g',n')$. 
%If $\pi(g,n)=\pi(g',n')$ then 
%$$
%(gP,gw_0P,gnw_0P)=(g'P,g'w_0P,g'n'w_0P). 
%$$
%As a consequence $g^{-1} g'$ lies in the stabilizer of $(P,w_0P)$ which means that $g^{-1}g' \in MA$. Without loss of generality, we can suppose that $g=e$ is the neutral element and $g'=\rho a$, where $\rho \in M$ and $a \in A$. In this way we obtain that 
%$$
%(P,w_0P,nw_0P)=(P,w_0P,\rho a n' a^{-1} \rho^{-1} w_0 P),
%$$
%where we exploited the fact that $w_0$ normalizes $MA$. By the uniqueness in the parametrization of Equation \eqref{eq:N:parametrization} we must have
%$$
%n=\rho a n' a^{-1} \rho^{-1}. 
%$$ 
%If we now plug the value of $n$ in the function $\varphi$ we get
%\begin{align*}
%\varphi(e,n)&=\Psi(n)K\\
%&=\Psi(\rho a n' a^{-1} \rho^{-1})K\\
%&=a\Psi(n')K\\
%&=a\Psi(n')\rho K\\
%&=\rho a \Psi(n') K=\varphi(\rho a , n'),
%\end{align*}
%where we used Equation \eqref{eq Psi M invariant} and \eqref{eq Psi A equivariant} to move from the second line to the third one and we concluded by exploiting the fact that $M<K$ commutes with $A$. Thus given two elements $(g,n)$ and $(g',n')$ such that $\pi(g,n)=\pi(g',n')$, we must have that $\varphi(g,n)=\varphi(g',n')$. Equivalently the map $\varphi$ factors through $\pi$. The induced map $\Phi$ is clearly $G$-equivariant and continuous by the universal property of quotients. This concludes the proof. 
%\end{proof}

\section{Generic tuples on the Furstenberg boundary and barycenter map}\label{sec barycenter}

Let $G$ be a connected semisimple Lie group with finite center. We consider a minimal parabolic subgroup $P<G$ and a maximal split torus $A<P$ with Lie algebra $\mathrm{Lie}(A)=\mathfrak{a}$. If $K$ denotes a maximal compact subgroup in $G$, recall that the \emph{Weyl group} $W$ is defined as the quotient 
$$
W:=N_K(\mathfrak{a})/Z_K(\mathfrak{a}),
$$
where $N_K$ and $Z_K$ are respectively the normalizer and the centralizer in $K$ of $\mathfrak{a}$ with respect to the adjoint representation. We fix a representative $w_0$ of the longest element in $W$. More generally, every time that we pick an element $w \in W$, we will tacitly assume to fix a representative of it. 

Two points $x,y \in G/P$ in the Furstenberg boundary are said to be \emph{opposite} if they lie in the same $G$-orbit as $(P,w_0P)$. Given two opposite points $x,y \in G/P$, we denote by $F_{x,y}$ the unique maximal flat determined by $x$ and $y$. In the particular case of $(P,w_0P)$, we denote it by $F_A$ and we call it the canonical maximal flat. The \emph{boundary} $\partial F \subset G/P$ of a maximal flat is the set of equivalence classes of Weyl chambers determined by $F$. For $F_A$ we have
$$
\partial F_A = \{ wP \ | \ w \in W\}.
$$
Recall that the $G$-action on the set of maximal flats is transitive and it holds that $\partial(gF_A)=g(\partial F_A)$. 

\begin{dfn}\cite{BucSavProj} \label{Def generic} Let $q\geq 1$. We define the set $(G/P)^{(q)}$ of \emph{generic} $q$-tuples in $(G/P)^q$ as follows: 
 \begin{itemize}
\item If $q=1$, take $(G/P)^{(1)}:=G/P$. 
\item If $q=2$, it is the set of opposite tuples. 
\item For $q\geq 3$, if $w_0$ acts as $-1$ on $\mathfrak{a}$, we define the set of generic points $(G/P)^{(q)}$ to be the set of pairwise opposite points. 
\item For $q\geq 3$, if $w_0$ does not act as $-1$ on $\mathfrak{a}$, we define the set of generic points $(G/P)^{(q)}$ to consists of $q$-tuples $(x_1,\dots,x_q)\in (G/P)^q$ such that $x_k$ is opposite to every point in the boundary of $\partial F_{x_i,x_j}$, for every distinct  $1 \leq i,j,k \leq q$.  
\end{itemize}
\end{dfn}

Observe that when the rank of $G$ is one, generic $q$-tuples are precisely $q$-tuples of pairwise distinct points. 

\begin{lem}\label{lem symmetrically weyl opposite}
Given a connected semisimple Lie group $G$ with finite center and a minimal parabolic subgroup $P<G$, the subset $(G/P)^{(q)}$ of generic $q$-tuples is open and dense and hence it has full measure. 
\end{lem}

\begin{proof}
Suppose that $w_0$ does not acts as $-1$ on $\mathfrak{a}=\mathrm{Lie}(A)$. We start showing that $(G/P)^{(3)}$ is an open dense subset of full measure. Given a point $x \in G/P$ we denote by 
$$
\mathrm{Opp}_x:=\{ y \in G/P \ | \ y \ \text{is opposite to} \ x \},
$$
and we define
$$
(G/P)_{\mathrm{opp}}:=\bigcap_{w \in W} \mathrm{Opp}_{wP}. 
$$
Notice that $(G/P)_{\mathrm{opp}}$ is open and dense in $G/P$, since each $\mathrm{Opp}_w$ is open and dense in $G/P$.  

The map
$$
\pi:G \times (G/P)_{\mathrm{opp}} \longrightarrow (G/P)^{(3)}, \ \ \pi(g,x):=(gP,gw_0P,gx),
$$
is an open surjection which parametrizes the subset of generic triples. Thus $(G/P)^{(3)}$ is open and dense in $(G/P)^3$. Finally $(G/P)^{(q)}$ is open and dense, being a finite intersection of open and dense subsets of $(G/P)^q$.

In the case where $w_0$ acts as $-1$ on $\mathfrak{a}$, the claim fowllows from the fact that opposite pairs $(G/P)^{(2)}$ form an open and dense subset in $(G/P)^2$. 
\end{proof}

The relevance of the subset $(G/P)^{(q)}$ of generic $q$-tuples relies on the existence of a \emph{barycenter map} onto the symmetric space $G/K$. More precisely, the authors \cite[Corollary 4]{BucSavProj} proved the existence of a $G$-equivariant symmetric continuous map
$$
\mathrm{bar}_q:(G/P)^{(q)} \longrightarrow G/K
$$
which we call \emph{barycenter}. We will exploit the map $\mathrm{bar}_3$ to obtain a vanishing result in the computation of the spectral sequence for continuous cochains. 

\section{Proof of Theorem \ref{thm_continuous}}\label{sec continuous theorem}

To prove Theorem \ref{thm_continuous} we adapt the spectral sequences from \cite{Monod} and \cite{BucSavAlt} to the continuous setting. Most of the adaptations are straightforward except for two crucial issues: 1) It is not clear that one of the spectral sequence associated to the double complex we consider vanishes, which is the reason that we cannot exploit that the other spectral sequence degenerates and conclude that there is an isomorphism with the  measurable cohomology. 2) The proof of the triviality of the columns $q\geq 3$ in the first page of the first spectral sequence by Monod \cite[Proposition 5.1]{Monod} does not work in the continuous setting. Our proof (Theorem \ref{thm vanishing gp unbounded}) will depend on the existence of  a $G$-equivariant symmetric continuous \emph{barycenter map} \cite{BucSavProj} 
\begin{equation}\label{eq intro bar 3}
\mathrm{bar}_3:(G/P)^{(3)} \rightarrow G/K.
\end{equation}
Here $G/K$ is the Riemannian globally symmetric space associated to $G$. When $G=\mathrm{Isom}^\circ(\mathbb{H}^n_\mathbb{R})$, the barycenter is precisely the barycenter of the ideal triangle defined by $x,y,z \in \partial \mathbb{H}^n_{\mathbb{R}}$. 

Recall that $G$ is a connected semisimple Lie group with finite center. Fix a
minimal parabolic subgroup $P<G$ and a maximal split torus $A<P$. Let $K$ be
a maximal compact subgroup and $M<A$ be the centralizer of $A$. Let $w_0$ be a
representative of the longest element in the Weyl group.

We want to replace the bicomplex
\begin{equation} \label{measurable alternating bicomplex}
M^{p,q}:=L^0(G^{p+1},L^0_{\alt}((G/P)^q))^G
\end{equation}
we introduced \cite{BucSavAlt} for measurable (hence the letter $M$) alterning cochains (which
was itself based on Monod's bicomplex \cite{Monod}), with an analogous bicomplex, but this time using continuous functions. More precisely, we define
\begin{equation}\label{bicomplex Kpq}
C^{p,q}:=C_c(G^{p+1},C_{c,\alt}((G/P)^{(q)}))^G,
\end{equation}
the space of $G$-equivariant continuous functions on $G^{p+1}$ with values in the space of continuous alternating functions on $(G/P)^{(q)}$. Alternation is well-defined on $(G/P)^{(q)}$ since the latter is invariant with respect to the action of the symmetric group $\mathrm{Sym}(q)$. The vertical differential 
$$
d^\uparrow:C^{p,q} \rightarrow C^{p+1,q}
$$
is the usual homogeneous differential on the $G$-variable, whereas the horizontal one 
$$
d^\rightarrow:C^{p,q} \rightarrow C^{p,q+1}
$$
is the homogeneous differential on $(G/P)^{(q)}$ weighted with the sign $(-1)^{p+1}$, so that $d^\rightarrow d^\uparrow=d^\uparrow d^\rightarrow$. 

Since the subset of generic $q$-tuples has full measure in $(G/P)^q$, we have a well-defined inclusion of continuous functions into (classes of) measurable ones 
$$
C^{p,q} \longrightarrow M^{p,q},
$$
which respects both the vertical differential and the horizontal one, leading to a map of bicomplexes. We will want to exploit this map and all the information about the bicomplex of measurables functions. We start by studying the first page of the spectral sequence
$$
E_1^{p,q}=(H^q(C^{\ast,p},d^\uparrow),d_1=d^\rightarrow). 
$$
The $(p,q)$-entry of this first page is given by
$$
E_1^{p,q}=H^q_c(G,C_{c,\alt}((G/P)^{(p)}). 
$$

Recall that the alternation map 
$$
\mathrm{Alt}_p:C_c((G/P)^{(p)}) \rightarrow C_c((G/P)^{(p)}),
$$
$$
\mathrm{Alt}_p(f)(g_1P,\ldots,g_pP)=\sum_{\sigma \in \mathrm{Sym}(p)} \mathrm{sgn}(\sigma)f(g_{\sigma(1)}P,\ldots,g_{\sigma(p)}P)
$$
induces a splitting
$$
C_c((G/P)^{(p)}) = C_{c,\alt}((G/P)^{(p)}) \oplus C_{c,\nalt}((G/P)^{(p)})
$$
into alternating and non-alternating functions, which further induces a splitting
\begin{equation}\label{eq coefficients splitting}
H^q_c(G,C_c((G/P)^{(p)})) \cong H_c^q(G,C_{c,\alt}((G/P)^{(p)})) \oplus H_c^q(G,C_{c,\nalt}((G/P)^{(p)}))
\end{equation}
on the corresponding cohomology groups. 
 
The first column ($p=0$) of the first page of our spectral sequence is clearly $E_1^{0,q}=H^q_c(G)$. Let us investigates the columns $p=1$ and $p=2$. Before doing that, we remind the reader of the existence of a natural action of $w_0$ on the cohomology $H^q_c(A)$ induced by the adjoint representation. In particular, we can consider the map 
$$
\Pi^q_{w_0}:H^q_c(A) \rightarrow H^q_c(A), \  \ \ \alpha \mapsto \alpha-\mathrm{Ad}(w_0)(\alpha). 
$$
We say that an element $\alpha \in H^q_c(A)$ is $w_0$-\emph{invariant} if it lies in the kernel of $\Pi^q_{w_0}$, whereas we say that $\alpha$ is $w_0$-\emph{equivariant} if it lies in the image of $\Pi^q_{w_0}$. We denote by 
$$
H^q_c(A)^{w_0}:=\Ker(\Pi^q_{w_0})
$$
the subspace of $w_0$-invariant classes. 

\begin{lem}\label{lem first two columns}
For every $q \geq 0$ we have that
$$
H^q_c(G,C_c(G/P)) \cong H^q_c(A),
$$
$$
H^q_c(G,C_{c,\alt}((G/P)^{(2)})) \cong \mathrm{Im}(\Pi^q_{w_0}).
$$
\end{lem} 

\begin{proof}
There exists an isomorphism of cocomplexes
\begin{equation}\label{iso shapiro P}
\Psi:C_c(G^{q+1})^P \longrightarrow C_c(G^{q+1},C_c(G/P))^G, \ \  (\Psi f)(g_0,\ldots,g_q)(gP):=f(g^{-1}g_0,\ldots,g^{-1}g_q).
\end{equation}
which induces the following isomorphism
$$
H^q_c(G,C_c(G/P)) \cong H_c^q(P),
$$
for every $q \geq 0$. Now it is sufficient to recall that the inclusion $A \rightarrow P$ induces an isomorphism in cohomology \cite[Proposition 3.1]{Monod}, namely
$$
H^q_c(P) \cong H^q_c(A),
$$
for every $q \geq 0$. 

Recall that the set $(G/P)^{(2)}$ of pairs of opposite points is an open dense subset of $(G/P)^2$ which can be identified with the quotient $G/MA$ via the map
$$
G/MA \rightarrow (G/P)^{(2)},  \ \ \ gMA \mapsto (gP,gw_0P).
$$
The existence of this $G$-equivariant diffeomorphism implies the following isomorphism
$$
H^q_c(G,C_c((G/P)^{(2)})) \cong H^q_c(G,C_c(G/MA)),
$$
for every $q \geq 0$. By substituting $P$ with $MA$ in Equation \eqref{iso shapiro P}, we obtain the isomorphism
$$
H^q_c(G,C_c(G/MA)) \cong H_c^q(MA),
$$
for every $q \geq 0$. Since $M$ is compact and centralizes $A$, we can apply \cite[Theorem 9.1]{Blanc} to get 
$$
H^q_c(MA) \cong H^q_c(A),
$$
for every $q \geq 0$. Moreover, the same of proof as \cite[Proposition 8]{BucSavAlt} shows that the splitting given by Equation \eqref{eq coefficients splitting} corresponds to the splitting of $H^q_c(A)$ into $w_0$-invariant and $w_0$-equivariant classes. More precisely, we have that 
$$
H^q_c(G,C_{c,\nalt}((G/P)^{(2)})) \cong H^q_c(A)^{w_0}, \ \ \ H^q_c(G,C_{\alt}((G/P)^{(2)})) \cong \mathrm{Im}(\Pi^q_{w_0}), 
$$
for every $q \geq 0$. This concludes the proof. 
\end{proof}

The next crucial step is to show triviality of the columns $E^{p,\ast}_1$ for $p\geq 3$. This is the first instance where a completely new ingredient is needed. The strategy will be to show  that the cohomology of the cocomplex $C_c(G^{*+1},C_c((G/P)^{(p)})^G$ is isomorphic to the cohomology of 
\begin{equation}\label{def Cpq}
C_K^{*,p}:=C_c((G/K)^{*+1},C_c((G/P)^{(p)}))^G.
\end{equation}
Then we will see that the latter cohomology group is trivial in the proof of Theorem \ref{thm vanishing gp unbounded}.

%\begin{dfn}\label{K invariant cochain}
%Let $\alpha \in C^{p,q}$. We say that $\alpha$ is a \emph{right-$K$-invariant} cochain if for every $k_0,\ldots,k_p \in K$ we have that 
%$$
%\alpha(g_0k_0,\ldots,g_pk_p)(h_1P,\ldots,h_qP)=\alpha(g_0,\ldots,g_p)(h_1P,\ldots,h_qP). 
%$$
%Denote by $C^{p,q}_K\subset $C^{p,q}$ the subset of right-$K$-invariant cochains. 
%\end{dfn}

\begin{lem}\label{K invariant representative}
The evaluation on the base point $K$, 
$$C_c((G/K)^{*+1},C_c((G/P)^{(p)}))^G\longrightarrow  C_c(G^{*+1},C_c((G/P)^{(p)}))^G$$
induces an isomorphism in cohomology
$$ H^q(C_c((G/K)^{*+1},C_c((G/P)^{(p)}))^G)\cong H^q( C_c(G^{*+1},C_c((G/P)^{(p)}))^G).$$
\end{lem}

\begin{proof}
Since $K$ is compact, it admits a finite Haar measure $\mu_K$. We define a left inverse of the evaluation on $K$ by 
$$
\alpha_K:C_c(G^{q+1},C_c((G/P)^{(p)}))^G\longrightarrow C_c((G/K)^{q+1},C_c((G/P)^{(p)}))^G$$
by
\begin{align*}
&\alpha_K(f)(g_0K,\dots,g_qK)(h_1P,\ldots ,h_pP)=\\
&\int_K \ldots \int_K f(g_0k_0,\ldots,g_qk_q)(h_1P,\ldots,h_pP)d\mu_K(k_0)\ldots d\mu_K(k_q), \end{align*}
for $f\in C_c(G^{q+1},C_c((G/P)^{(p)}))^G$. 
%Given a $G$-invariant cocycle, we can explicitly construct a cohomologous right-$K$-invariant cochain by setting
%$$
%\alpha_K:G^{p+1} \rightarrow C_c((G/P)^{(q)}),
%$$
%$$
%\alpha_K(g_0,\ldots,g_p)(h_1P,\ldots,h_qP):=
%$$
Notice that the integrals are all well-defined since $\mu_K$ is finite and $f$ is continuous. It is obvious that $\alpha_K$ is a left inverse to the evaluation on $K$. 

The composition of  $\alpha_K$ followed by evaluation is simply
\begin{align*}
&\alpha(f)(g_0,\ldots,g_q) (h_1P,\ldots ,h_pP)=\\
&\int_K \ldots \int_K f(g_0k_0,\ldots,g_qk_q)(h_1P,\ldots,h_pP)d\mu_K(k_0)\ldots d\mu_K(k_q), \end{align*}
for $f\in C_c(G^{q+1},C_c((G/P)^{(p)}))^G$. 

To prove that $\alpha$ is homotopic to the identity we define the following chain homotopy
$$
H:C_c(G^{q+1},C_c((G/P)^{(p)}))^G \rightarrow C_c(G^{q},C_c((G/P)^{(p)}))^G,
$$
$$
(H f)(g_0,\ldots,g_{q-1})(h_1P,\ldots,h_pP):=
$$
$$
\sum_{i=0}^{q-1} (-1)^i \int_K \ldots \int_K f(g_0,\ldots,g_i,g_ik_i,\ldots,g_{q-1}k_{q-1})(h_1P,\ldots,h_pP)d\mu_K(k_i)\ldots d\mu_K(k_{q-1}). 
$$
A standard computation shows that 
$$
\mathrm{Id}-\alpha=d(H\alpha),
$$
and the statement is proved. 
\end{proof}

\begin{thm}\label{thm vanishing gp unbounded}
For every $q \geq 1$ and $p \geq 3$ we have that 
$$
H_c^q(G,C_c((G/P)^{(p)})) \cong 0.
$$
\end{thm}

\begin{proof}
Thanks to Lemma \ref{K invariant representative} it is sufficient to construct a contracting homotopy for the cocomplex 
$$
(C_K^{p,q},d^\uparrow).
$$

Since in the coefficient module we have at least three points on the boundary, we will exploit the barycenter map of Section \ref{sec barycenter}. More precisely, we define
$$
H:C_K^{p,q} \rightarrow C_K^{p-1,q},
$$ 
$$
(H f)(g_0K,\dots,g_{p-1}K)(h_1P,\ldots,h_qP):=
$$
$$
f(\mathrm{bar}_3(h_1P,h_2P,h_3P),g_0K,\ldots,g_{p-1}K)(h_1P,\ldots,h_qP). 
$$
We claim that $H$ is well-defined. For any $f \in C_K^{p,q}$, the function $Hf$ is continuous, being the composition of continuous functions. In a similar way, $Hf$ is $G$-invariant by the $G$-equivariance of the map $\mathrm{bar}_3$. A direct computation shows that 
$$
f=H(df)+d(Hf),
$$
which precisely means that $H$ is a contracting homotopy for the cocomplex, as claimed. This concludes the proof. 
\end{proof}

\begin{cor}\label{cor vanishing 3 points}
For every $q \geq 1$ and every $p \geq 3$ we have that 
$$
H^q_c(G,C_{c,\alt}((G/P)^{(p)})) \cong 0. 
$$
\end{cor}

\begin{proof}
By Equation \eqref{eq coefficients splitting} we have a splitting 
$$
H^q_c(G,C_c((G/P)^{(p)})) \cong H_c^q(G,C_{c,\alt}((G/P)^{(p)})) \oplus H_c^q(G,C_{c,\nalt}((G/P)^{(p)})).
$$
Theorem \ref{thm vanishing gp unbounded} implies that the left-hand side of the previous equation vanishes. As a consequence the two summands on the right-hand side must vanish as well, and the statement is proved. 
\end{proof}

\begin{proof}[Proof of Theorem \ref{thm_continuous}]

We start by computing explicitly the first page of the spectral sequence
$$
E_1^{p,q}=(H^q(C^{\ast,p},d^\uparrow),d_1=d^\rightarrow). 
$$
By Corollary \ref{cor vanishing 3 points} we have that $E_1^{p,q} \cong 0$ whenever $q \geq 1$ and $p \geq 3$. The first column $E_1^{0,q}$ is naturally isomorphic to the continuous cohomology of $G$, namely $H^q_c(G)$. By Lemma \ref{lem first two columns} we have that the columns with coefficients having either one or two points on the boundary are isomorphic to 
$$
E^{1,q}_1 \cong H^q_c(A), \ \ \ E^{2,q}_1 \cong \mathrm{Im}(\Pi^q_{w_0}),
$$
for $q \geq 1$.
Finally the bottom row coincides with the $G$-invariant continuous alternating functions on $(G/P)^{(p)}$, that is
$$
E^{p,0}_1 \cong C_{c,\alt}((G/P)^{(p)})^G.
$$ 

\begin{figure}[!h]
\centering
\begin{tikzpicture}
  \matrix (m) [matrix of math nodes,
             nodes in empty cells,
             nodes={minimum width=9ex,
                    minimum height=9ex,
                    outer sep=-3pt},
             column sep=2ex, row sep=-1ex,
             text centered,anchor=center]{
         q      &         &          &          & \\
          \cdots & \cdots & \cdots & \cdots   & \cdots \\
          3    & \ H^3_c(G) \  & \ H^3_c(A) \ &\  \mathrm{Im}(\Pi^3_{w_0}) \ & 0 & \ \cdots & \\
          2    &  \ H^2_c(G) \ &\  H^2_c(A) \  & \ \mathrm{Im}(\Pi^2_{w_0}) \ & 0 & \  \cdots &  \\
          1    & \ H^1_c(G) \ & \ H^1_c(A) \ & \  \mathrm{Im}(\Pi^1_{w_0}) \  & 0 & \   \cdots & \\
          0     & \ \mathbb{R}\   &\  \C_{c}(G/P)^G \ & \  C_{c,\alt}((G/P)^{(2)})^G  \ & \   C_{c,\alt}((G/P)^{(3)})^G \  &  \ \cdots & \\
    \quad\strut &   0  &  1  &  2  &  3  &  \cdots & p \strut \\};

%lines are counted starting from the top 

%arrows in the 3rd line
\draw[->](m-3-2.east) -- (m-3-3.west)node[midway,above ]{$0$};
\draw[->](m-3-3.east) -- (m-3-4.west)node[midway,above ]{$\Pi^3_{w_0}$};
\draw[->](m-3-4.east) -- (m-3-5.west);

%arrows in the 2nd line
\draw[->](m-4-2.east) -- (m-4-3.west)node[midway,above ]{$0$};
\draw[->](m-4-3.east) -- (m-4-4.west)node[midway,above ]{$\Pi^2_{w_0}$};
\draw[->](m-4-4.east) -- (m-4-5.west);

%arrows in the 1st line
\draw[->](m-5-2.east) -- (m-5-3.west)node[midway,above ]{$0$};
\draw[->](m-5-3.east) -- (m-5-4.west)node[midway,above ]{$\Pi^1_{w_0}$};
\draw[->](m-5-4.east) -- (m-5-5.west);

%arrows in the 1st line
\draw[->](m-6-2.east) -- (m-6-3.west)node[midway,above ]{$\delta$};
\draw[->](m-6-3.east) -- (m-6-4.west)node[midway,above ]{$\delta$};
\draw[->](m-6-4.east) -- (m-6-5.west)node[midway,above ]{$\delta$};



\draw[thick] (m-1-1.east) -- (m-7-1.east) ;
\draw[thick] (m-7-1.north) -- (m-7-7.north) ;
\end{tikzpicture}
\caption{The first page $E_1$}\label{fig first page unbounded}
\end{figure}

The differential $d_1=d^\rightarrow$ on the first row boils down to the usual homogeneous differential $\delta$. For $q \geq 1$, the differential 
$$
d_1:H^q_c(G) \cong E^{0,q}_1 \longrightarrow E^{1,q}_1 \cong H^q_c(A)
$$
coincides with the restriction map multiplied by $(-1)^{p+1}$. By \cite[Corollary 3]{Wienhard} the restriction map vanishes identically, thus the same holds for the differentials from the first column to the second one. Finally, the same proof as \cite[Theorem 3]{BucSavAlt} shows that the differentials from the second column to the third one are conjugated to the operator $\Pi^q_{w_0}$. We depict the first page $E_1$ in Figure \ref{fig first page unbounded}.

The second page $E_2$ is now easily computed. The isomorphism $H^0_c(G \curvearrowright G/P) \cong \mathbb{R}$ implies that $E^{0,0}_2 \cong E^{0,1}_2 \cong 0$. For $p \geq 2$, on the row $q=0$ appears the cohomology on $(G/P)^{(p)}$ shifted by one, that is 
$$
E^{p,0}_2 \cong H^{p-1}_c(G \curvearrowright G/P). 
$$
The continuous cohomology of $G$ on the first column $p=0$ is preserved, namely $E^{0,q}_2 \cong H^q_c(G)$. By the surjectivity of the differentials from the column $q=1$ to the column $q=2$, we obtain that 
$$
E^{1,q}_2 \cong H^q_c(A)^{w_0}, \ \ \ E^{2,q}_2 \cong 0. 
$$
We report the second page $E_2$ in Figure \ref{fig second page unbounded}. 

\begin{figure}[!h]
\centering
\begin{tikzpicture}
  \matrix (m) [matrix of math nodes,
             nodes in empty cells,
             nodes={minimum width=9ex,
                    minimum height=9ex,
                    outer sep=-3pt},
             column sep=2ex, row sep=-1ex,
             text centered,anchor=center]{
	q      &         &          &          & \\
          \cdots & \cdots & \cdots & \cdots   & \cdots \\
          3    & \ H^3_c(G) \  & \ H^3_c(A)^{w_0} \ &\  0 \ & 0 & \ \cdots & \\
          2    &  \ H^2_c(G) \ &\  H^2_c(A)^{w_0} \  & \  0 \ & 0 & \  \cdots &  \\
          1    & \ H^1_c(G) \ & \ H^1_c(A)^{w_0} \ & \  0 \  & 0 & \   \cdots & \\
          0     & \ 0\   &\  0 \ & \  H^1_{c,\alt}(G \curvearrowright G/P)  \ & \  H^2_{c,\alt}(G \curvearrowright G/P) \  &  \ \cdots & \\
    \quad\strut &   0  &  1  &  2  &  3  &  \cdots & p \strut \\};

%lines are counted starting from the top 

%arrows in the 3rd line
\draw[->](m-3-2) -- (m-4-4);
\draw[->](m-3-3) -- (m-4-5);

%arrows in the 2nd line
\draw[->](m-4-2) -- (m-5-4);
\draw[->](m-4-3) -- (m-5-5);

%arrows in the 1st line
\draw[->](m-5-2) -- (m-6-4);
\draw[->](m-5-3) -- (m-6-5);

\draw[thick] (m-1-1.east) -- (m-7-1.east) ;
\draw[thick] (m-7-1.north) -- (m-7-7.north) ;
\end{tikzpicture}
\caption{The second page $E_2$} \label{fig second page unbounded}
\end{figure}

Since we do not know if the spectral sequence degenerates or not, we will exploit now the map between bicomplexes
\begin{equation} \label{eq bicomplexes}
C^{p,q} \longrightarrow M^{p,q},
\end{equation}
where $M^{p,q}$ is the bicomplex of Equation \eqref{measurable alternating bicomplex}. We denote by
$$
E^{p,q}_{1,m}:=(H^q(M^{\ast,p},d^\uparrow),d_1=d^\rightarrow).
$$
the first page of the spectral sequence generated by the bicomplex on the right-hand side of Equation \eqref{eq bicomplexes}. The page $E_{1,m}$ is reported in \cite[Figure 3]{BucSavAlt}. The maps determined by Equation \eqref{eq bicomplexes} induce maps
\begin{equation}\label{eq first pages}
E_1^{p,q} \longrightarrow E_{1,m}^{p,q}
\end{equation}
for every $p,q \geq 0$. We claim that those maps are isomorphism on each column for $q \geq 1$. When $p=0$ this is precisely the statement of Austin and Moore \cite[Theorem A]{AM}. Let us prove that the map $E^{1,q}_1 \rightarrow E^{1,q}_{1,m}$ induced up to conjugation by the inclusion 
\begin{equation}\label{eq inclusion continuous}
C_c(G^{q+1})^P \longrightarrow L^0(G^{q+1})^P,
\end{equation}
is an isomorphism. Notice that the inclusion of Equation \eqref{eq inclusion continuous} can be restricted to $A$ invariant cochains obtaining the following commutative diagram
$$
\xymatrix{
C_c(G^{q+1})^P \ar[rr] \ar[d] && L^0(G^{q+1})^P \ar[d] \\
C_c(G^{q+1})^A \ar[rr] && L^0(G^{q+1})^A.
}
$$
By \cite[Proposition 3.1]{Monod} the columns of the previous diagram are isomorphisms in cohomology. Additionally, the inclusion on the bottom row can be decomposed as
$$
C_c(G^{q+1})^A \rightarrow C_c(A^{q+1})^A \rightarrow L^0(A^{q+1})^A \rightarrow L^0(G^{q+1})^A,
$$
where the first function is the usual restriction map, the second one is the inclusion and the third one is induced by the $A$-equivariant projection $\pi_A:G \rightarrow A$, where $\pi_A$ is defined in terms of the Iwasawa decomposition $G=ANK$ as $\pi_A(ank)=a$, for $a\in A,\ n\in N$ and $k\in K$.  By \cite[Theorem A]{AM} and \cite[Theorem 2]{Moore} all of these maps induce isomorphisms in cohomology, thus the claim follows. The same argument also applies to the column $p=2$, so 
$$
E^{2,q}_1 \longrightarrow E^{2,q}_{1,m}
$$
are all isomorphisms, again for $q \geq 1$. 

We now pass to the second pages, where the page $E_{2,m}$ is depicted in \cite[Figure 4]{BucSavAlt}. The existence of the maps in  \eqref{eq first pages} implies that there exists a natural map 
\begin{equation}\label{eq second pages}
E^{p,q}_2 \longrightarrow E^{p,q}_{2,m},
\end{equation}
for every $p, q \geq 0$. Again for $q \geq 1$ these maps are all isomorphisms by \cite[Theorem 1]{AM}. As a consequence we obtain a commutative diagram
\begin{equation}\label{diagram A}
\xymatrix{
H^q_c(A)^{w_0} \ar[rr] \ar[rd] && H^q_{c,\alt}(G \curvearrowright G/P) \ar[ld]\\
& H^q_{m,\alt}(G \curvearrowright G/P), &
}
\end{equation}
where $H^q_{m,\alt}(G \curvearrowright G/P)$ is the cohomology of the cocomplex $(L^0_{\alt}((G/P)^{\ast+1})^G,\delta)$. By \cite[Theorem 3]{BucSavAlt} the map 
$$
H^q_c(A)^{w_0} \longrightarrow H^q_{m,\alt}(G \curvearrowright G/P)
$$
is injective, so the map 
$$
H^q_c(A)^{w_0} \longrightarrow H^q_{c,\alt}(G \curvearrowright G/P)
$$
is injective as well by the commutativity of Diagram \eqref{diagram A}. If we now pass to the quotients, the maps given by \eqref{eq second pages} determine another commutative diagram
\begin{equation}\label{eq diagram G}
\xymatrix{
H^q_c(G) \ar[rr] \ar[rd] && H^q_{c,\alt}(G \curvearrowright G/P)/H^q_c(A)^{w_0} \ar[ld] \\
& H^q_{m,\alt}(G \curvearrowright G/P)/H^q_c(A)^{w_0}. &
}
\end{equation}
Since \cite[Theorem 3]{BucSavAlt} guarantees that 
$$
H^q_c(G) \longrightarrow H^q_{m,\alt}(G \curvearrowright G/P)/H^q_c(A)^{w_0}
$$
is an isomorphism, we must have that 
$$
H^q_c(G) \longrightarrow H^q_{c,\alt}(G \curvearrowright G/P)/H^q_c(A)^{w_0}
$$
is injective and that 
$$
H^q_{c,\alt}(G \curvearrowright G/P)/H^p_c(A)^{w_0} \longrightarrow H^q_{m,\alt}(G \curvearrowright G/P)/H^p_c(A)^{w_0}
$$
is surjective. In particular 
$$
H^q_{c,\alt}(G \curvearrowright G/P)\longrightarrow H^q_{m,\alt}(G \curvearrowright G/P)
$$
also is surjective. 
\end{proof}




\section{Proof of Theorem \ref{thm_continuous_bounded}} \label{sec bounded case}
We follow the same strategy as for Theorem \ref{thm_continuous} in the previous section. We construct a bicomplex analogous to the one in  \eqref{bicomplex Kpq}, but this time using continuous bounded functions. More precisely, we define
$$
C_{b}^{p,q}:=C_{c,b}(G^{p+1},C_{c,b,\alt}(((G/P)^{(q)}))^G, 
$$
where $(G/P)^{(q)}$ is the space of generic $q$-tuples. The vertical differential 
$$
d^\uparrow:C^{p,q}_b \longrightarrow C^{p+1,q}_b
$$
is simply the homogeneous differential on the $G$-variable, whereas the horizontal one
$$
d^\rightarrow:C^{p,q}_b \longrightarrow C^{p,q+1}_b
$$
is the homogeneous differential on $(G/P)^{(q)}$ weighted with the sign $(-1)^{p+1}$. 

As before, the bicomplex $(C_b^{p,q},d^\rightarrow, d^\uparrow)$ determines two different spectral sequences. In contrast to the unbounded setting, the fact that we are now dealing with bounded functions, allows us to prove vanishing of the  first spectral sequence
\begin{equation}\label{eq first spectral bounded}
^I_bE^{p,q}_1:=(H^q(C^{p,\ast}_b,d^{\rightarrow}),d_1=d^\uparrow). 
\end{equation}

\begin{prop}\label{first bounded zero}
Let $^I_bE^{p,q}_1$ be the first page of the spectral sequence defined by \eqref{eq first spectral bounded}. Then the spectral sequence degenerates immediately, that is $^I_bE^{p,q}_1=0$ for every $p,q \geq 0$. 
\end{prop}

\begin{proof}
The first page of the spectral sequence is obtained by considering the horizontal differential $d^\rightarrow$, namely it is obtained by the cohomology of the cocomplex
\begin{equation}\label{eq cocomplex horizontal}
\rightarrow C_{c,b}(G^{p+1},C_{c,b,\alt}(((G/P)^{(q-1)}))^G \rightarrow C_{c,b}(G^{p+1},C_{c,b,\alt}(((G/P)^{(q)}))^G \rightarrow .
\end{equation}
By \cite[Proposition 7.4.12]{Monodbook} the cohomology of the above cocomplex is the same as the one of its \emph{inhomogeneous} variant, which is obtained by getting rid of $G$-invariance and by deleting one $G$-variable. 
\begin{equation}\label{eq inhomogeneous bounded}
\rightarrow C_{c,b}(G^{p},C_{c,b,\alt}(((G/P)^{(q-1)})) \rightarrow C_{c,b}(G^{p},C_{c,b,\alt}(((G/P)^{(q)})) \rightarrow .
\end{equation}
It is worth noticing that the differential on the boundary variable is still the homogeneous one, suitably weighted with a sign. By the proof of \cite[Lemma 7.5.5]{Monodbook} the cocomplex
$$
0 \rightarrow \mathbb{R} \rightarrow  C_{c,b}(G/P) \rightarrow C_{c,b,\alt}((G/P)^{(2)}) \rightarrow
$$
admits a contracting homotopy given by integration along the first variable and hence it is exact. To obtain the cohomology of the cocomplex of Equation \eqref{eq cocomplex horizontal} we are applying the functor $C_{c,b}(G^p, \ \cdot \ )$, which is exact by \cite[Lemma 8.2.4]{Monodbook}. This concludes the proof. 
\end{proof} 

We now turn to the second spectral sequence, namely the one with first page
$$
^{II}_bE_1^{p,q}:=(H^q(C_b^{\ast,p},d^\uparrow),d_1=d^\rightarrow). 
$$
The $(p,q)$-entry of this first page is given by
\begin{equation}\label{pq entry bounded}
^{II}_bE_1^{p,q}:=H_{c,b}^q(G,C_{c,b,\alt}((G/P)^{(p)})).
\end{equation}
It is worth noticing that the restriction of the alternation map $\mathrm{Alt}_p$ to $C_{c,b}((G/P)^{(p)})$ is continuous with respect to the supremum norm. As a consequence we have a splitting
\begin{equation}\label{splitting bounded}
C_{c,b}((G/P)^{(p)}) \cong C_{c,b,\alt}((G/P)^{(p)}) \oplus C_{c,b,\nalt}((G/P)^{(p)})
\end{equation}
and the decomposition holds at the level of Banach spaces (in particular each subspace is closed and complemented). Thus, by \cite[Corollary 8.2.10]{Monodbook} we have a similar splitting at cohomology level
\begin{equation}\label{splitting bounded cohomology}
H^q_{c,b}(G,C_{c,b}((G/P)^{(p)})) \cong H^q_{c,b}(G,C_{c,b,\alt}((G/P)^{(p)})) \oplus H^q_{c,b}(G,C_{c,b,\nalt}((G/P)^{(p)})). 
\end{equation}

Our aim is to identify all the terms in the page $^{II}_bE_1$. We are going to prove that all the entries which are not either on the first column or on the first row must vanish. We start proving this statement for the coefficient module with one or two points in the boundary. We will exploit the following version of Eckmann-Shapiro induction.

\begin{prop}\label{continuous induction iso}
Let $G$ be a connected semisimple Lie group with finite center. Let $L < G$ be a closed subgroup. Then we have the following isomorphism
$$
H^q_{c,b}(G,C_{c,b}(G/L)) \cong H^q_{c,b}(L),
$$
for every $q \geq 0$. 
\end{prop}

\begin{proof}
Before starting the proof, we recall that, given a Banach $G$-module $E$, its \emph{maximal continuous submodule} is defined by
$$
\mathcal{C}E:=\{ v \in E \ | \ g.v \rightarrow v \ \textup{when} \  g \rightarrow e\}. 
$$
In the computation of continuous bounded cohomology, all the information is contained in the maximal continuous submodule. More precisely, Monod \cite[Proposition 6.1.5]{Monodbook} proved the following isometric isomorphism
$$
H^q_{c,b}(G,E) \cong H^q_{c,b}(G,\mathcal{C}E),
$$
for any Banach $G$-module $E$ and any degree $p \geq 0$. As a consequence, we have that 
$$
H^q_{c,b}(G,C_{c,b}(G/L)) \cong H^q_{c,b}(G,\mathcal{C}C_{c,b}(G/L)),
$$
for any $q \geq 0$.

By the universal property of quotients, the Banach space $C_{c,b}(G/L)$ is isomorphic to $L$-invariant continuous bounded functions on $G$ (notice that $L$-invariants are considered with respect to the right action of $G$ on the bimodule $C_{c,b}(G)$). In this way we obtain that 
$$
H^q_{c,b}(G,\mathcal{C}C_{cb}(G/L)) \cong H^q_{c,b}(G,\mathcal{C}(C_{c,b}(G)^L)) \cong H^q_{c,b}(G,\mathcal{C}C_{c,b}(G)^L),
$$
where the last isomorphism follows by the fact that $\mathcal{C}( \ \cdot \ )$ and the $L$-invariants commute. As shown in the proof of \cite[Proposition 4.4.2]{Monodbook}, we have the following isomorphism
$$
\mathcal{C}C_{c,b}(G) \cong \mathcal{C}L^\infty(G). 
$$
As a consequence, we obtain that 
$$
H^q_{c,b}(G,\mathcal{C}C_{c,b}(G)^L) \cong H^q_{cb}(G,\mathcal{C}L^\infty(G)^L) \cong H^q_{c,b}(G,\mathcal{C}L^\infty(G/L)), 
$$
where the last isomorphism follows from the fact that 
$$
L^\infty(G)^L \cong L^\infty(G/L),
$$
in virtue of Fubini's theorem. 

Thanks to the standard induction isomorphism \cite[Proposition 10.1.3]{Monodbook} in bounded cohomology, we finally obtain that 
$$
H^q_{c,b}(G,\mathcal{C}L^\infty(G/L)) \cong H^q_{c,b}(G,L^\infty(G/L)) \cong H^q_{c,b}(L).
$$
This concludes the proof. 
\end{proof}

\begin{cor}\label{cor second third columns}
For every $q \geq 1$ we have that 
$$
H^q_{c,b}(G,C_{c,b}(G/P)) \cong H^q_{c,b}(G,C_{c,b,\alt}((G/P)^{(2)})) \cong 0. 
$$
\end{cor}

\begin{proof}
By Proposition \ref{continuous induction iso} we can write 
$$
H^q_{c,b}(G,C_{c,b}(G/P)) \cong H^p_{c,b}(P) \cong 0,
$$
where the latter is zero because of the amenability of $P$ \cite[Corollary 7.5.11]{Monodbook}. 


By  \eqref{splitting bounded cohomology}, the space $H^q_{c,b}(G,C_{c,b,\alt}((G/P)^{(2)}))$ is a direct summand of the cohomology group 
$$
H^p_{c,b}(G,C_{cb}((G/P)^{(2)})) \cong H^p_{cb}(G,C_{cb}(G/MA)).
$$
Exploiting once again Proposition \ref{continuous induction iso}, we immediately obtain that
$$
 H^q_{c,b}(G,C_{c,b}(G/MA)) \cong H^q_{c,b}(MA) \cong 0,
$$
where the vanishing statement follows by the amenability of $MA$ \cite[Corollary 7.5.11]{Monodbook}. Thus $H^q_{c,b}(G,C_{c,b,\alt}((G/P)^{(2)}))$ must vanish as well, and the statement is proved.
\end{proof}

We are left to show that, for $q \geq 1$, all the remaining columns but the first one are trivial.

\begin{thm}\label{thm vanishing gp bounded}
For every $q \geq 1$ and $p \geq 3$ we have that 
$$
H^q_{c,b}(G,C_{c,b}((G/P)^{(p)})) \cong 0. 
$$
\end{thm} 

\begin{proof}
The strategy is the same as Theorem \ref{thm vanishing gp unbounded}. First, it is worth noticing that the same argument that we exploited to show Lemma \ref{K invariant representative} can be adapted in the bounded context to prove that the cocomplex 
$$
(_bC_K^{p,q},d^\uparrow):=((C_{c,b}((G/K)^{p+1},C_{c,b}((G/P)^{(q)}))^G,d^\uparrow),
$$
computes the cohomology group $H^q_{c,b}(G,C_{c,b}((G/P)^{(p)}))$. As a consequence it is sufficient to build a contracting homotopy for the cocomplex $_bC_K^{p,q}$. The same homotopy defined in the proof of Theorem \ref{thm vanishing gp unbounded} works fine, since it preserves boundedness. This concludes the proof. 
\end{proof}

\begin{cor}\label{cor bounded 3 points}
For every $q \geq 1$ and $p \geq 3$ we have that
$$
H^q_{c,b}(G,C_{c,b,\alt}((G/P)^{(p)})) \cong 0.
$$
\end{cor}

\begin{proof}
By  \eqref{splitting bounded cohomology}, the cohomology group $H^q_{c,b}(G,C_{c,b,\alt}((G/P)^{(p)}))$ is a direct summand of $H^q_{c,b}(G,C_{c,b}((G/P)^{(p)}))$, which vanishes by Theorem \ref{thm vanishing gp bounded}.
\end{proof}

\begin{proof}[Proof of Theorem \ref{thm_continuous_bounded}]
We consider the first page of the spectral sequence 
$$
^{II}_bE^{p,q}_1:=(H^q(C^{\ast,p}_b,d^\uparrow),d_1:=d^\rightarrow)
$$
This page is particularly easy: in fact by Corollary \ref{cor second third columns} and Corollary \ref{cor bounded 3 points} we have that $^{II}_bE^{p,q}_1 \cong 0$ for every $p \geq 1$ and $q \geq 1$. We are left only with the first column and the first row. On the first column appears the bounded cohomology of $G$, namely
$^{II}_bE^{0,q}_1 \cong H_{c,b}^q(G)$. In a similar way, on the bottom row we get back continuous bounded alternating functions on $(G/P)^{(p)}$, that is
$^{II}_bE^{p,0}_1 \cong C_{c,b,\alt}((G/P)^{(p)})$. Additionally the differential $d_1=d^\rightarrow$ boils down to the usual homogeneous differential $\delta$ on $(G/P)^{(p)}$. We report the first page $^{II}_bE^{p,q}_1$ in Figure \ref{first page bounded}.

\begin{figure}[!h]
\centering
\begin{tikzpicture}
  \matrix (m) [matrix of math nodes,
             nodes in empty cells,
             nodes={minimum width=9ex,
                    minimum height=9ex,
                    outer sep=-3pt},
             column sep=2ex, row sep=-1ex,
             text centered,anchor=center]{
         q      &         &          &          & \\
          \cdots & \cdots & \cdots & \cdots   & \cdots \\
          3    & \ H^3_{c,b}(G) \  & \ 0 \ &\  0 \ & 0 & \ \cdots & \\
          2    &  \ H^2_{c,b}(G) \ &\  0 \  & \ 0 \ & 0 & \  \cdots &  \\
          1    & \ H^1_{c,b}(G) \ & \ 0 \ & \  0 \  & 0 & \   \cdots & \\
          0     & \ \mathbb{R}\   &\  \C_{c,b}(G/P)^G \ & \  C_{c,b,\alt}((G/P)^{(2)})^G  \ & \   C_{c,b,\alt}((G/P)^{(3)})^G \  &  \ \cdots & \\
    \quad\strut &   0  &  1  &  2  &  3  &  \cdots & p \strut \\};

%lines are counted starting from the top 

%arrows in the 3rd line
\draw[->](m-3-2.east) -- (m-3-3.west)node[midway,above ]{$0$};
\draw[->](m-3-3.east) -- (m-3-4.west)node[midway,above ]{$0$};
\draw[->](m-3-4.east) -- (m-3-5.west);

%arrows in the 2nd line
\draw[->](m-4-2.east) -- (m-4-3.west)node[midway,above ]{$0$};
\draw[->](m-4-3.east) -- (m-4-4.west)node[midway,above ]{$0$};
\draw[->](m-4-4.east) -- (m-4-5.west);

%arrows in the 1st line
\draw[->](m-5-2.east) -- (m-5-3.west)node[midway,above ]{$0$};
\draw[->](m-5-3.east) -- (m-5-4.west)node[midway,above ]{$0$};
\draw[->](m-5-4.east) -- (m-5-5.west);

%arrows in the 1st line
\draw[->](m-6-2.east) -- (m-6-3.west)node[midway,above ]{$\delta$};
\draw[->](m-6-3.east) -- (m-6-4.west)node[midway,above ]{$\delta$};
\draw[->](m-6-4.east) -- (m-6-5.west)node[midway,above ]{$\delta$};



\draw[thick] (m-1-1.east) -- (m-7-1.east) ;
\draw[thick] (m-7-1.north) -- (m-7-7.north) ;
\end{tikzpicture}
\caption{The first page $^{II}_bE_1$}\label{first page bounded}
\end{figure}

The second page $^{II}_bE^{p,q}_2$ is simple as well: the only thing to modify is the bottom row, where we now have the cohomology on the boundary shifted by $1$, namely 
$$
^{II}_bE^{p,0}_2 \cong H_{c,b}^{p-1}(G \curvearrowright G/P). 
$$
We report the second page $^{II}_bE^{p,q}_2$ in Figure \ref{second page bounded}. 


\begin{figure}[!h]
\centering
\begin{tikzpicture}
  \matrix (m) [matrix of math nodes,
             nodes in empty cells,
             nodes={minimum width=9ex,
                    minimum height=9ex,
                    outer sep=-3pt},
             column sep=2ex, row sep=-1ex,
             text centered,anchor=center]{
         q      &         &          &          & \\
          \cdots & \cdots & \cdots & \cdots   & \cdots \\
          3    & \ H^3_{c,b}(G) \  & \ 0 \ &\  0 \ & 0 & \ \cdots & \\
          2    &  \ H^2_{c,b}(G) \ &\  0 \  & \  0 \ & 0 & \  \cdots &  \\
          1    & \ H^1_{c,b}(G) \ & \ 0 \ & \  0 \  & 0 & \   \cdots & \\
          0     & \ 0\   &\  0 \ & \  H^1_{c,b,\alt}(G \curvearrowright G/P)  \ & \  H^2_{c,b,\alt}(G \curvearrowright G/P) \  &  \ \cdots & \\
    \quad\strut &   0  &  1  &  2  &  3  &  \cdots & p \strut \\};

%lines are counted starting from the top 

%arrows in the 3rd line
\draw[->](m-3-2) -- (m-4-4);
\draw[->](m-3-3) -- (m-4-5);

%arrows in the 2nd line
\draw[->](m-4-2) -- (m-5-4);
\draw[->](m-4-3) -- (m-5-5);

%arrows in the 1st line
\draw[->](m-5-2) -- (m-6-4);
\draw[->](m-5-3) -- (m-6-5);

\draw[thick] (m-1-1.east) -- (m-7-1.east) ;
\draw[thick] (m-7-1.north) -- (m-7-7.north) ;
\end{tikzpicture}
\caption{The second page $^{II}_bE_2$} \label{second page bounded}
\end{figure}
 
Since the spectral sequence $^I_bE_1$ degenerates immediately by Proposition \ref{first bounded zero}, by \cite[Appendice A]{Guichardet} the spectral sequence $^{II}_bE_1$ must converge to the same limit. The only way to obtain this convergence is that
$$
d_{q+1}:H^q_{c,b}(G) \longrightarrow H^q_{c,b,\alt}(G \curvearrowright G/P)
$$
is an isomorphism. 

Following the proof of Theorem \ref{thm_continuous}, we can consider the bounded analogue of the bicomplex $M^{p,q}$, namely
$$
M^{p,q}_b=L^\infty(G^{p+1},L^\infty_{\mathrm{alt}}((G/P)^q))^G,
$$
with the same vertical and horizontal differentials. We clearly have an inclusion
$$
C^{p,q}_b \longrightarrow M^{p,q}_b,
$$
which is a map of bicomplexes. We denote by
$$
^{II}_bE^{p,q}_{1,m}:=(H^q(M^{\ast,p}_b,d^\uparrow),d_1=d^\rightarrow)
$$
the first page of the second spectral sequence generated by the measurable bicomplex. The page $^{II}_bE_{1,m}$ is the same as the one reported in Figure \ref{first page bounded} up to substituing continuous cochains with measurable ones. The natural maps 
$$
^{II}_bE^{p,q}_{1,m} \rightarrow ^{II}_bE^{p,q}_{1}
$$
are isomorphism for every $p \geq 0$ and $q \geq 1$ by \cite[Proposition 7.5.1]{Monodbook}. By passing to the second pages, for every $q \geq 0$, we obtain the following diagram
$$
\xymatrix{
H^q_{c,b}(G) \ar[rr]^{d_{q+1,m}} \ar[rd]^{d_{q+1}} && H^q_{m,b,\alt}(G \curvearrowright G/P) \ar[ld]\\
& H^q_{c,b,\alt}(G \curvearrowright G/P), &
}
$$
where $H^q_{m,b,\mathrm{alt}}(G \curvearrowright G/P)$ is the cohomology of the cocomplex $(L^\infty((G/P)^{\ast+1})^G,\delta)$ and $d_{q+1,m}$ is the differential of order $p+1$ in the measurable bicomplex. Since we proved that $d_{q+1}$ is an isomorphism and $d_{q+1,m}$ is an isomorphism by \cite{Monod}, the statement follows.  
\end{proof}


\section{Proof of Proposition 3}


A direct consequence of the proofs of Theorem \ref{thm_continuous} and \ref{thm_continuous_bounded} is that we have the following commutative diagram
\begin{equation}\label{diagram comparison}
\xymatrix{
H^q_{cb}(G) \ar[d]_{\mathrm{comp}^q_G} \ar@{^{(}->>}[rr]^{\hspace{-40pt} d_\infty}_\cong && H^q_{cb,\alt}(G \curvearrowright G/P) \ar[d]^{p\circ \mathrm{comp}^q_{G \curvearrowright G/P}}\\
H^q_c(G) \ar@{^{(}->}[rr]^{\hspace{-40pt} d_\infty} && H^q_{c,\alt}(G \curvearrowright G/P)/H^{k-1}_c(A)^{w_0},
}
\end{equation}
where $p:H^q_{c,\alt}(G \curvearrowright G/P)\rightarrow H^q_{c,\alt}(G \curvearrowright G/P)/H^{q-1}_c(A)^{w_0}$ is the natural projection and we tacitly identify $H^{q-1}_c(A)^{w_0}$ with the image $d_\infty(H^{q-1}_c(A)^{w_0})\subset H^q_{c,\alt}(G \curvearrowright G/P)$. The functions $d_\infty$ appearing in the top and bottom arrows are determined by the differentials of the associated spectral sequences and  are both injective. 

Since the bottom $d_\infty$ is injective, the injectivity of $\mathrm{comp}^q_G$ is equivalent to the injectivity of the composition $d_\infty \circ  \mathrm{comp}^q_G= p\circ \mathrm{comp}^q_{G \curvearrowright G/P}\circ d_\infty$. Now the top $d_\infty$ is furthermore known to be an isomorphism so the latter map is injective if and only if $p\circ \mathrm{comp}^q_{G \curvearrowright G/P}$ is. The latter claim is equivalent to $\mathrm{comp}^q_{G \curvearrowright G/P}$ being injective, and for the kernel of $p$ to intersect the image of $\mathrm{comp}^q_{G \curvearrowright G/P}$ trivially, which is precisely the statement of the Proposition. 

\section{Injectivity of the comparison map}\label{sec injectivity}

We are finally ready to prove the main injectivity results of the paper.

\begin{proof}[Proof of Theorem \ref{thm boundary deg 4}] We first deal with the case $n=2$. We identify the boundary of the hyperbolic plane with the real projective line $\mathbb{P}^1(\mathbb{R})$. Recall that in the real rank one case generic tuples in $\mathbb{P}^1(\mathbb{R})^{(k)}$ simply are distinct tuples.

Consider a $\mathrm{PSL}(2,\mathbb{R})$-invariant continuous bounded alternating cocycle
$$
b:\mathbb{P}^1(\mathbb{R})^{(5)} \longrightarrow \mathbb{R}. 
$$
Suppose that there exists a $\mathrm{PSL}(2,\mathbb{R})$-invariant continuous alternating function 
$$
f:\mathbb{P}^1(\mathbb{R})^{(4)} \longrightarrow \mathbb{R}
$$
such that $\delta f=b$. We will show that $f$ is bounded, which implies the desired injectivity. We define
$$\begin{array}{rrcl}
F:&\mathbb{P}^1(\mathbb{R}) \setminus \{\infty,0,1\}& \longrightarrow& \mathbb{R},\\
&x&\longmapsto &f(\infty,0,1,x).
\end{array}$$
We will actually show that $F$ is bounded, which readily implies that $f$ is bounded as well by the transitivity of $\mathrm{PSL}(2,\mathbb{R})$ on positively oriented triples and the fact that $f$ is alternating. Recall that for distinct $x_i\in P^1(\mathbb{R})$, by the $\mathrm{PSL}(2,\mathbb{R})$-invariance of $f$, we have that 
$$f(x_0,x_1,x_2,x_3)=f(\infty,0,1,[x_0,x_1,x_2,x_3])=F([x_0,x_1,x_2,x_3]),$$
 where 
$$[x_0,x_1,x_2,x_3]=\frac{x_0-x_2}{x_0-x_3}\cdot \frac{x_1-x_3}{x_1-x_2}$$
denotes the usual cross ratio. 

By the continuity of $F$, we only need to prove its boundedness in some neighborhood of $\infty$, $0$ and $1$. In fact exploiting again the fact that $f$ is alternating it will suffice to prove boundedness of $F$ on some interval $(1-\delta,1)$, for any $\delta>0$. Let us see how this implies boundedness of $F$ on the three neighborhoods: For $x \in (1,1+\delta']$ it is sufficient to notice that 
\begin{equation}\label{x to 1/x}
F(x)=f(\infty,0,1,x)=-f(\infty,0,x,1)=-f\left(\infty,0,1,\frac{1}{x}\right)=-F\left(\frac{1}{x}\right)
\end{equation}
where we exploited the fact that $f$ is alternating. The latter term is bounded since $1/x<1$ lies in a neighborhood of $1$.  The function $F$ is thus uniformly bounded around $1$ (where the latter value is excluded). For the boundedness around $0$ we have
$$F(x)=f(\infty,0,1,x)=-f(\infty,1,0,x)=-F(1-x),$$
so that the boundedness around $0$ is equivalent to the boundedness around $1$. Finally for the boundedness around $\infty$ we use the boundedness around $0$ and the relation (\ref{x to 1/x})


%We can exploit the fact that $f$ is alternating twice to get the following relations
%\begin{equation}\label{F alternation}
%F(x)=F\left( 1- \frac{1}{x} \right), \ \ F(x)=F\left(\frac{1}{1-x}\right).
%\end{equation}

Let $0<\delta<1$. We will show that $F$ is bounded on $(1-\delta,1)$. 
We apply the relation $\delta f =b$ to a $5$-tuple $(\infty,0,1,x,y)$ of distinct points
to obtain
\begin{equation}\label{F cocycle}
F(x)-F(y)+F\left(\frac{y}{x}\right)-F\left( \frac{1-y}{1-x} \right) + F\left( \frac{x(1-y)}{y(1-x)} \right) \sim 0,
\end{equation}
where the symbol $\sim$ means that the left-hand side is uniformly at bounded distance from the right-hand side. Indeed the left-hand side is equal to $b(\infty,0,1,x,y)$, which is bounded by our assumption on $b$. 

We now plug into  Equation \eqref{F cocycle} the value $y=x^2$ (which is allowed since for $x\in \mathbb{R}\setminus \{0,1\}$ the $5$-tuple $(\infty,0,1,x,x^2)$ consists of distinct points) to obtain
\begin{equation}\label{eq x minus x}
2F(x)-F(x^2)-F(1+x)+F\left(\frac{1+x}{x}\right) \sim 0.
\end{equation}
If $x$ lies in a neighborhood of $1$, then $1+x$ and $(1+x)/x$ both lie in a neighborhood of $2$ where $F$ is continuous and hence bounded. In particular, for $x\in [1-\delta,1)$, the expression \eqref{eq x minus x} rewrites as 
\begin{equation}\label{eq F twice}
F(x^2) \sim 2F(x),
\end{equation}
More precisely, there must exist  $C>0$ such that for every $x \in [1-\delta,1)$ it holds that 
\begin{equation}\label{eq C}
\frac{F(x^2)}{2}-C \leq F(x) \leq \frac{F(x^2)}{2}+C. 
\end{equation}
Since 
$$
[1-\delta,1)=\bigcup_{k \geq 0} [(1-\delta)^{\frac{1}{2^k}},(1-\delta)^{\frac{1}{2^{k+1}}}),
$$
there exists some $k\geq 0$ such that $x \in [(1-\delta)^{\frac{1}{2^k}},(1-\delta)^{\frac{1}{2^{k+1}}})$. By \eqref{eq C} we have  
$$
\frac{F(x^2)}{2}-C \leq F(x) \leq \frac{F(x^2)}{2}+C,
$$
with $x^2 \in [(1-\delta)^{\frac{1}{2^{k-1}}},(1-\delta)^{\frac{1}{2^k}})$. If $k> 0$ we can apply \eqref{eq C} again to obtain
$$
\frac{F(x^4)}{4}-C\left(1+\frac{1}{2}\right) \leq F(x) \leq \frac{F(x^4)}{4}+C\left(1+\frac{1}{2}\right),
$$
with $x^4 \in [(1-\delta)^{\frac{1}{2^{k-2}}},(1-\delta)^{\frac{1}{2^{k-1}}})$. Finally we iterate this procedure $k+1$ times, until we obtain 
$$
\frac{F(x^{2^{k+1}})}{2^{k+1}}-C\left(1+\frac{1}{2}+\ldots+\frac{1}{2^k}\right) \leq F(x) \leq \frac{F(x^{2^{k+1}})}{2^{k+1}}+C\left(1+\frac{1}{2}+\ldots+\frac{1}{2^k}\right),
$$
with $x^{2^k} \in [(1-\delta)^2,1-\delta)$. The continuity of $F$ on the closure of $[(1-\delta)^2,1-\delta)$ ensures that $F(x^{2^k})$ is bounded. Hence the convergence of the geometric series guarantees that $F(x)$ is bounded for $x \in [1-\delta,1)$ and concludes the proof for $n=2$. 



When $n=3$ the argument is similar:  We identify the Furstenberg boundary with the complex projective line $\mathbb{P}^1(\mathbb{C})$. As before, we suppose that a $\mathrm{PSL}(2,\mathbb{C})$-invariant continuous bounded alternating cocycle 
$$
c:\mathbb{P}^1(\mathbb{C})^{(5)} \longrightarrow \mathbb{R}
$$
can be written as $\delta f=c$, where $f$ is a $G$-invariant continuous alternating function
$$
f:\mathbb{P}^1(\mathbb{C})^{(4)} \longrightarrow \mathbb{R}.
$$
We define 
$$\begin{array}{rrcl}
F: &\mathbb{P}^1(\mathbb{C}) \setminus \{\infty,0,1 \} &\longrightarrow& \mathbb{R}\\
&x&\longmapsto &f(\infty,0,1,x),
\end{array}$$
it is sufficient to prove the boundedness of $F$. As in the real case, it will be enough to show the boundedness of $F$ on a neighborhood of $1$. The same equations as in the $n=2$ case imply the boundedness of $F$ around $0$ and $\infty$ as well. Furthermore by (\ref{x to 1/x}) is enough to prove boundedness in the intersection of a neighborhood of $1$ and the closed unit disc centered at $1$. 

Fix some small $\delta>0$ and let $U$ be  given as
$$U=\{z\in \mathbb{C}\setminus \{1\} \mid 1-\delta <|z|\leq 1, \ -\delta < \mathrm{arg}(z)<\delta\}.$$
Precisely as in the $n=2$ case, for $x\in U$ we have
$$
2F(x)-F(x^2) \sim 0.
$$
Thus applying the same inductive estimates  it is enough to observe that for every $z\in U$ there exists $k$ such that $z^{2^k}$ belongs to the closure of 
$$\{ z\in \mathbb{C}\setminus \{1\} \mid (1-\delta)^2<|z|\leq 1, -2\delta < \mathrm{arg}(z) < 2\delta\} \setminus U,$$
on which $F$ is bounded by continuity. This concludes the proof for $n=3$.

The case $n \geq 4$ follows immediately from the case $n=3$: Let 
$$
b:(\partial \mathbb{H}^n_{\mathbb{R}})^{(5)} \longrightarrow \mathbb{R}
$$
be a $G$-invariant continuous bounded alternating cocycle and suppose that it can be written as $\delta f=b$, where $f$ is a $G$-invariant continuous alternating function
$$
f: (\partial \mathbb{H}^n_{\mathbb{R}})^{(4)} \longrightarrow \mathbb{R}. 
$$
Fix a copy of $\partial \mathbb{H}^3_{\mathbb{R}} \subset \partial \mathbb{H}^n_{\mathbb{R}}$. By the case $n=3$, we know that the restriction $f|_{(\partial \mathbb{H}^3_{\mathbb{R}})^{(4)}}$ is bounded. This implies that $f$ is bounded as well: an arbitrary $4$-tuple $(x_0,x_1,x_2,x_3)$ of boundary points lies in an $\mathrm{Isom}^+(\mathbb{H}^n)$-translate of $\partial \mathbb{H}^3$, so there exists $g \in \mathrm{Isom}^\circ(\mathbb{H}^n_{\mathbb{R}})$ such that $gx_0,gx_1,gx_2,gx_3 \in \partial \mathbb{H}^3_{\mathbb{R}}$, and the claim follows by the $G$-invariance of $f$. This concludes the proof.
\end{proof}


\bibliographystyle{alpha}
\begin{thebibliography}{20}

\bibitem{AM} Austin T., Moore C. C. 
\emph{Continuity properties of measurable group cohomology}, Math. Ann. {\bf 356} (2013), n. 3, 885--937

\bibitem{Blanc} Blanc P. 
\emph{Sur la cohomologie continue des groupes localement compacts}, Ann. Sci. \'Ecole Norm. Sup.{\bf 12} (1979), n.2,137--168.

\bibitem{Bloch} Bloch S. J. \emph{Higher regulators, algebraic K-theory and zeta functions of elliptic curves}, CRM Monograph Series 11, Amer. Math. Soc. (2000).

\bibitem{BBI18} Bucher M., Burger M. and Iozzi A., \emph{The bounded borel class and complex representations of
3-manifold groups}, Duke Math. {\bf 167} (2018), no. 17, 3129--3169.

\bibitem{BucSavExp} Bucher M, Savini A., \emph{Some explicit cocycles on the Furstenberg boundary for products of isometries of hyperbolic spaces and $\mathrm{SL}(3,\mathbb{K})$}, preprint.

\bibitem{BucSavAlt} Bucher M., Savini A., \emph{Alternating cochains on Furstenberg boundaries and measurable cohomology}, preprint.

\bibitem{BucSavProj} Bucher M., Savini A., \emph{Projections from Furstenberg boundaries onto maximal flats and barycenter maps}, preprint. 

\bibitem{BM99} Burger M., Monod N, \emph{Bounded cohomology of lattices in higher rank Lie groups}, J. Eur.
Math. Soc. \textbf{1} (1999), 199--235.

\bibitem{BuMo} Burger M., Monod. N., \emph{Continuous bounded cohomology and applications to rigidity
theory}, Geometric and Functional Analysis {\bf 12} (2002), n. 219--280.

\bibitem{BuMo2} Burger M., Monod N., \emph{On and around the bounded cohomology of SL2}, from: Rigidity in dynamics and geometry, {M. Burger, A. Iozzi editors}, Springer, Berlin (2002), 19--37.

\bibitem{DLC} De la Cruz Mengual C., \emph{The degree-three bounded cohomology of complex Lie groups of classical type},
preprint.

\bibitem{Guichardet}  Guichardet A., \emph{Cohomologie des groupes topologiques et des alg\'ebres de Lie}, Textes
Math\'ematiques, vol. 2, CEDIC, Paris, 1980.

\bibitem{HO} Hartnick T. and Ott A., \emph{Bounded cohomology via partial differential equations}. Geom. Topol.
{\bf 19} (6), (2015), 3603--3643.

\bibitem{Helgason} Helgason S., \emph{Differential geometry, Lie groups, and symmetric spaces}, Corrected reprint
of the 1978 original, Grad. Stud. Math. 34, Amer. Math. Soc., Providence, R.I., 2001.

\bibitem{Monodbook} Monod N., \emph{Continuous bounded cohomology of locally compact groups}. Number 1758 
Lecture notes in Mathematics. Springer-Verlag, Berlin, 2001.

\bibitem{MonodSLn} Monod N., \emph{Stabilization for SLn in bounded cohomology}, Proceedings of the first JAMS
Symposium, Matsushima 2002. Contemp. Math. {\bf 347} (2004), 191--202.

\bibitem{MonodICM} Monod N., \emph{An invitation to bounded cohomology}, Proceedings of the ICM 2006, Volume II, 1183--1211.

\bibitem{Monod} Monod N., \emph{The cohomology of semi-simple Lie groups, viewed from infinity}, Trans. Amer. Math. Soc. B {\bf 9} (2022), 144--159 

\bibitem{Moore} Moore C. C., \emph{Groups extensions and cohomology for locally compact groups, III}, Trans. Amer. Math. Soc. {\bf 221}
(1976), n. 1, 1--33

\bibitem{Pieters} Pieters H. \emph{The boundary model for the continuous cohomology of $\mathrm{Isom}^+(\mathbb{H}^n)$}, Groups, Geometry and Dynamics {\bf 12} (2018), n. 4, 1239--1263.

\bibitem{Wienhard} Wienhard A.,
\emph{Remarks on and around bounded differential forms}, Pure Appl. Math. Q. {\bf 8} (2012), n. 2, 479--496

%\bibitem{Zimmer} Zimmer R. J., \emph{Ergodic theory and semisimple Lie groups}, Monographs in Mathematics, vol. {\bf 81}, Birkauser, Verlag, Basel (1984)

%%%%%%%%%%%%%%%%%




\end{thebibliography}

\end{document}




